\documentclass[10pt,a4paper]{article}
\usepackage{amsmath,amssymb,amsthm,graphicx}
\usepackage{hyperref}
\usepackage{fontspec}
\title{Fine-grained randomized simulation of unique-witness computation by hitting-set reconstruction}
\author{Halley Goldberg, Valentine Kabanets, Zhenjian Lu, Igor C. Oliveira}
\date{doi:10.4230/LIPIcs.CCC.2022.16; CC BY 4.0}
\begin{document}
\maketitle
\noindent\textit{Editorial scope.} Counted original Theorem67 (Theorem2 Item1) once for its independent long unique-witness hitting-set/avoider reconstruction mechanism. Quasilinear randomized average-case easiness and original time bounds remain essential assumptions. Complete AppendixB and necessary full earlier proof chains included. Fine-grained AM Theorem45 is uncounted context after limited comparison with067; support lemmas, parameter/time variants and direct corollaries are not separate tasks. Original mathematical statements and author-labelled proof adaptations are retained without added proofs or corrections.\par\medskip
\noindent\textit{Allowed standard prerequisites (editorial).}\begin{itemize}
\item Imported short-output black-box hitting-set generator and reconstruction [10, Theorem4.3], DPG reconstruction [11], standard quasilinear classes and error amplification. Exact imported hitting-set statement70 retained. Complete source-local approximation, weak symmetry, depth, Claim71, deterministic-avoider construction and witness reconstruction proof retained. Author fine-grained adaptations47/49/50 explicitly reference full earlier proofs, all of which are included here.
\end{itemize}\medskip
% Faithful official CCC2022 article16 excerpts: prerequisites and complete pKt proof chain for the UP main theorem.
\section*{2 Preliminaries (necessary author context)}
\subsection*{2.1 Notation}
We say that a function $p:\mathbb N\to\mathbb N$ is non-decreasing if $p(\ell)\geq\ell$ for every integer $\ell\geq0$. Throughout the paper, we will rely on constructions whose associated complexity parameters (e.g., running time) can be bounded by some polynomial $p$. We will implicitly assume that the polynomials provided in these results are non decreasing in order to avoid including this condition in every statement. This can be done without loss of generality by the monotonicity of our bounds.
For a probability distribution $D_n$, we use $\operatorname{supp}(D_n)$ to denote its support. Given an element $x$, we let $D_n(x)$ denote the probability of $x$ under $D_n$. We denote by $U_n$ the uniform distribution over $n$-bit strings.
Let $D_1$ and $D_2$ be probability distributions supported over a set $X$, $A:X\to\{0,1\}$, and $\varepsilon\geq0$. We say that $A$ $\varepsilon$-distinguishes $D_1$ and $D_2$ if
\[\left|\Pr_{x\sim D_1}[A(x)=1]-\Pr_{x\sim D_2}[A(x)=1]\right|\geq\varepsilon.\]
If $A$ is a randomized algorithm, each probability in the expression above also takes into account the internal randomness of $A$, which will be denoted by $r_A$.
\subsection*{2.2 Average-case complexity}
Recall that a pair $(L,\mathcal D)$ is a distributional problem if $L\subseteq\{0,1\}^*$ and $\mathcal D=\{D_n\}_{n\geq1}$ is an ensemble of probability distributions, where each $D_n$ is supported over $\{0,1\}^*$.
Let $\mathcal D=\{D_n\}_{n\geq1}$ be an ensemble of distributions. We say that $\mathcal D\in\mathrm{PSamp}$ if there is a randomized polytime algorithm $A$ such that, for every $n\geq1$, $A(1^n)$ is distributed according to $D_n$. More generally, we use $D_n\in\mathrm{Samp}[T(n)]/a(n)$ to denote that $D_n$ can be sampled by an algorithm that runs in time $T(n)$ and has advice complexity $a(n)$.
We let $\mathrm{DistNP}$ denote the set of distributional problems $(L,\mathcal D)$ with $L\in\mathrm{NP}$ and $\mathcal D\in\mathrm{PSamp}$.
\paragraph{Definition 4 ($\mathrm{Avg}_\delta\mathrm{BPP}$ [3]).} Let $(L,\mathcal D)$ be a distributional problem, and $\delta:\mathbb N\to[0,1]$. We say that $(L,\mathcal D)\in\mathrm{Avg}_\delta\mathrm{BPP}$ if there is randomized polytime algorithm $A$ such that
\begin{enumerate}
\item for every $n>0$, and every $x\in\operatorname{supp}(D_n)$,
\[\Pr_{r_A}[A(x;n)\notin\{L(x),\bot\}]\leq1/4,\]
\item for every $n>0$,
\[\Pr_{x\sim D_n}\left[\Pr_{r_A}[A(x;n)=\bot]\geq1/4\right]\leq\delta(n),\]
\end{enumerate}
where $r_A$ denotes the internal randomness of $A$. Such an algorithm $A$ is called an randomized errorless heuristic for $(L,\mathcal D)$ with failure probability at most $\delta$.
\paragraph{Definition 5 (AvgBPP [3]).}\footnote{Some authors define $\mathrm{AvgBPP}$ as $\bigcap_{c>0}\mathrm{Avg}_{n^{-c}}\mathrm{BPP}$. All our results also hold with respect to that definition.} We say that $(L,\mathcal D)$ is in $\mathrm{AvgBPP}$ if there exist a randomized algorithm $A$ and a polynomial $p$ such that
\begin{enumerate}
\item for every $n,\delta>0$, and every $x\in\operatorname{supp}(D_n)$,
\[\Pr_{r_A}[A(x;n,\delta)\notin\{L(x),\bot\}]\leq1/4,\]
\item for every $n,\delta>0$,
\[\Pr_{x\sim D_n}\left[\Pr_{r_A}[A(x;n,\delta)=\bot]\geq1/4\right]\leq\delta(n),\]
\item for every $n,\delta>0$, and every $x\in\operatorname{supp}(D_n)$, $A(x;n,\delta)$ runs in time at most $p(n/\delta)$.
\end{enumerate}
Such an algorithm $A$ is called a randomized errorless heuristic scheme for $(L,\mathcal D)$.
\paragraph{Lemma 6.} Let $(L,\mathcal D)$ be a distributional problem in $\mathrm{Avg}_{n^{-1}}\mathrm{BPP}$. There exists a randomized polynomial-time algorithm $B$ such that
\begin{enumerate}
\item for every $n>0$, and every $x\in\operatorname{supp}(D_n)$, $B(x;n)\in\{0,1\}$,
\item for every $n>0$, and every $x\in\operatorname{supp}(D_n)$, if $x\in L$, then $\Pr_{r_B}[B(x;n)=1]\geq1-1/n$, and
\item for every $n>0$, $\Pr_{x\sim D_n,r_B}[B(x;n)=L(x)]\geq1-3/n$.
\end{enumerate}
\paragraph{Proof.} Let $A$ be an $\mathrm{Avg}_{n^{-1}}\mathrm{BPP}$ algorithm for $(L,\mathcal D)$ as in Definition 4. Let $A'$ be the algorithm obtained, using standard amplification techniques,\footnote{In more detail, $A'$ can be obtained from $A$ by running it $O(n)$ times and selecting the most common output.} such that for all $x\in\operatorname{supp}(D_n)$,
\[\Pr_{r_{A'}}[A'(x;n)\notin\{L(x),\bot\}]\leq1/n\tag{1}\]
and
\[\Pr_{x\sim D_n}\left[\Pr_{r_{A'}}[A'(x;n)=\bot]\geq1/n\right]\leq1/n.\tag{2}\]
Let $B$ be the randomized algorithm which, on input $x\in\operatorname{supp}(D_n)$, simulates $A'(x;n)$ and outputs $1$ if $A'(x;n)\in\{1,\bot\}$ or $0$ if $A'(x;n)=0$. By Equation (1), if $x\in L$, then
\[\Pr_{r_B}[B(x;n)=0]=\Pr_{r_{A'}}[A'(x;n)\notin\{L(x),\bot\}]\leq1/n.\]
Let $\mathrm{bad}:=\{x\in\operatorname{supp}(D_n)\mid\Pr_{r_{A'}}[A'(x;n)=\bot]\geq1/n\}$. By Equation (2), at most a $(1/n)$-measure of $x\in\operatorname{supp}(D_n)$ are in bad. Thus,
\begin{align*}
\Pr_{x\sim D_n,r_B}[B(x;n)\neq L(x)]
&\leq\Pr_{x,r_{A'}}[A'(x;n)=\bot]+\Pr_{x,r_{A'}}[A'(x;n)\notin\{L(x),\bot\}]\\
&\leq\Pr_x[x\in\mathrm{bad}]+\Pr_{x,r_{A'}}[A'(x;n)=\bot\mid x\notin\mathrm{bad}]\\
&\qquad+\Pr_{x,r_{A'}}[A'(x;n)\notin\{L(x),\bot\}]\\
&\leq1/n+1/n+1/n=3/n.
\end{align*}
This completes the proof.\hfill$\triangleleft$
\subsection*{2.4 Direct product generators (DPGs)}
\paragraph{Definition 13 (Direct Product Generator (DPG) [11, Definiton 3.10]).} For $k,n\in\mathbb N$, we define the $k$-wise direct product generator to be the function
\[\mathrm{DP}_k:\{0,1\}^n\times\{0,1\}^{nk}\to\{0,1\}^{nk+k}\]
such that
\[\mathrm{DP}_k(x;z^1,\ldots,z^n):=(z^1,\ldots,z^k,x\cdot z^1,\ldots,x\cdot z^k).\]
\paragraph{Lemma 14 (DPG Reconstruction [11, Lemma 3.14]).} For any $n,k\in\mathbb N$ with $k\leq2n$, and $\varepsilon>0$, there exists a pair of algorithms $A$ and $\mathrm{Recon}^{(-)}$ such that
\begin{itemize}
\item $\mathrm{Recon}^D$ takes oracle access to an oracle $D:\{0,1\}^{nk+k}\to\{0,1\}$.
\item $A:\{0,1\}^n\times\{0,1\}^r\to\{0,1\}^k$ is called the advice function and is computable in time $\operatorname{poly}(n/\varepsilon)$.
\item $\mathrm{Recon}^{(-)}:\{0,1\}^k\times\{0,1\}^r\to\{0,1\}^n$ is called a reconstruction procedure and is computable in time $\operatorname{poly}(n/\varepsilon)$.\footnote{We note that if the oracle is a uniform algorithm that runs in time $t(m)$ on inputs of length $m$, the reconstruction procedure can be computed in time $t(nk+k)\cdot\operatorname{poly}(n/\varepsilon)$.}
\item The randomness complexity $r$ is at most $\operatorname{poly}(n/\varepsilon)$.
\item For any string $x\in\{0,1\}^n$ and any function $D$ that $\varepsilon$-distinguishes $\mathrm{DP}_k(x;U_{nk})$ from $U_{nk+k}$, it holds that
\[\Pr_{w\sim\{0,1\}^r}[\mathrm{Recon}^D(A(x,w),w)=x]\geq1/\operatorname{poly}(n/\varepsilon).\]
\end{itemize}
\subsection*{2.7 Encoding tuples (author context)}
We will also need to consider distributions that generate input instances that are tuples, while the only known information to the distributions is the length of the instances. To do this, we will encode the format as the length of the instances. More specifically, we use the following way to encode tuples [3].
\paragraph{Proposition 16 ([3]; see also [6, Proposition 3.1]).} There is a pair of polytime encoding-decoding algorithms $\mathrm{Enc}:\mathbb N^*\to\mathbb N$ and $\mathrm{Dec}:\mathbb N\to\mathbb N^*\cup\{\bot\}$ such that
\begin{enumerate}
\item For $k\in\mathbb N$ and $k$ integers $n_1,n_2,\ldots,n_k$, we write $\langle n_1,n_2,\ldots,n_k\rangle:=\mathrm{Enc}(n_1,n_2,\ldots,n_k)$, and we have
\begin{enumerate}
\item $\sum_{i=1}^k(2\cdot(\lfloor\log n_i\rfloor+2)+n_i)\leq\langle n_1,\ldots,n_k\rangle\leq O_k(\prod_{i=1}^k(n_i\cdot\log^2n_i))$.
\item $\mathrm{Dec}(\langle n_1,n_2,\ldots,n_k\rangle)=n_1,\ldots,n_k$. Here, Dec takes the binary representation of an integer and outputs $k$ integers which are also represented in binary.
\end{enumerate}
\item $\mathrm{Dec}(u)=\bot$ if $u\neq\langle\vec n\rangle$ for any $\vec n\in\mathbb N^*$.
\end{enumerate}
\paragraph{Proof.} Given $k$ integers $n_1,n_2,\ldots,n_k$, the encoding algorithm Enc first obtains the binary string
\[z:=\alpha_1\circ01\circ\alpha_2\circ01\circ\alpha_k\circ01\circ\operatorname{bin}(n_1)\circ\operatorname{bin}(n_2)\circ\cdots\circ\operatorname{bin}(n_k),\]
where for each $i\in[k]$, $\alpha_i$ is obtained from $\operatorname{bin}(|\operatorname{bin}(n_i)|)$ with every bit doubled. Then Enc outputs the integer $n$ whose binary representation is $z$. It is easy to see that given the binary representation of $n$, which is simply the string $z$, one can recover the $n_1,n_2,\ldots,n_k$ in time polynomial in $|z|$. Also, item (a) can be verified by a simple calculation.\hfill$\triangleleft$
\section*{3 Probabilistic Time-Bounded Kolmogorov Complexity (author excerpts)}
In this section, we formally define probabilistic Kolmogorov complexity and prove some useful properties of this notion, which will be used later in proving our main results.
\paragraph{Definition 17 ($\mathrm{pK}^t$).} For strings $x,y\in\{0,1\}^*$, a time bound $t\in\mathbb N$, an oracle $A$, and $\delta\in[0,1]$, the $\delta$-probabilistic $A$-oracle $t$-bounded Kolmogorov complexity of $x$ given $y$ is defined as
\[\mathrm{pK}_\delta^{t,A}(x\mid y):=\min\left\{k\ \middle|\ \Pr_{w\sim\{0,1\}^t}[\exists M\in\{0,1\}^k,M^A(w,y)\text{ outputs }x\text{ within }t\text{ steps}]\geq\delta\right\},\]
where $M$ is a RAM-machine.\footnote{We use the RAM model in this definition for convenience: it simplifies the time bound estimates for a machine $M$ that takes as input both randomness $w$ and an ``auxiliary'' string $y$, and allows both $w$ and $y$ to be as long as the running time of the RAM-machine $M$; cf. [29] for a similar definition in the case of deterministic time-bounded conditional Kolmogorov complexity. Alternatively, we could use a TM model, with inputs $w$ and $y$ presented on two different tapes. However, we prefer the RAM model as it's also useful in the context of fine-grained worst-case to average-case reductions [6].} We omit the subscript $\delta$ if $\delta=2/3$, omit the superscript $A$ if $A=\emptyset$, omit ``$\mid y$'' if $y$ is the empty string, and omit the superscript $t$ if $t=\infty$.
For simplicity, in the rest of this section we only consider $\mathrm{pK}^t$ without oracles. It is easy to see that all the results also hold with any oracle.
\paragraph{Lemma 21 (Success Amplification).} For any string $x\in\{0,1\}^n$, time bound $t\in\mathbb N$, and $0\leq\alpha<\beta\leq1$, we have
\[\mathrm{pK}_\beta^{O(qt/\alpha)}(x)\leq\mathrm{pK}_\alpha^t+O(\log(q/\alpha)),\]
where $q=\ln(1/(1-\beta))$.
\paragraph{Proof.} Suppose $\mathrm{pK}_\alpha^t(x)=k$. Then by definition, we have
\[\Pr_{w\sim\{0,1\}^t}[\exists M_w\in\{0,1\}^k\text{ such that }M_w(w)\text{ outputs }x\text{ within }t\text{ steps}]\geq\alpha.\]
Call such a $w$ good. After sampling $\ell$ independent $w^{(1)},w^{(2)},\ldots,w^{(\ell)}\in\{0,1\}^t$, the probability that no $w^{(i)}$ is good is at most $(1-\alpha)^\ell\leq e^{-\alpha\ell}$, which can be made at most $1-\beta$ by choosing $\ell=q/\alpha$. It follows that, with probability at least $\beta$ over a random $w=(w^{(1)},w^{(2)},\ldots,w^{(q/\alpha)})\in\{0,1\}^{qt/\alpha}$, there exists an index $1\leq i\leq q/\alpha$ (described with at most $\lceil\log(q/\alpha)\rceil$ bits) and a program $M_{w^{(i)}}\in\{0,1\}^k$ such that $M_{w^{(i)}}(w^{(i)})$ outputs $x$ within $O(qt/\alpha)$ steps. Hence, $\mathrm{pK}_\beta^{O(qt/\alpha)}(x)\leq k+O(\log(q/\alpha))$.\hfill$\triangleleft$
Note that the parameter $\alpha$ affects the time complexity in the resulting $\mathrm{pK}_\beta^{O(qt/\alpha)}$ bound. Lemma 21 is sufficient in applications where $\alpha\geq1/\operatorname{poly}(n)$.
\subsection*{3.2 Bounding $\mathrm{pK}^t$ via DPG reconstruction (selected original case)}
The following is a key lemma for $\mathrm{pK}^t$ that we will use in proving our main results related to randomized average-case complexity.
\paragraph{Lemma 22 ($\mathrm{pK}^t$ Reconstruction Lemma).} For $\varepsilon>0$, $x\in\{0,1\}^n$, $s\in\mathbb N$, and $k\in\mathbb N$ satisfying $k\leq2n$, let $D$ be a randomized algorithm that takes an advice string $\beta$ and runs in time $t_D$ such that $D$ $\varepsilon$-distinguishes $\mathrm{DP}_k(x;U_{nk})$ from $U_{nk+k}$. Then there is a polynomial $p_{\mathrm{DP}}$ such that
\[\mathrm{pK}^{\widetilde O(t_D)\cdot p_{\mathrm{DP}}(n/\varepsilon)}(x\mid\beta)\leq k+\log p_{\mathrm{DP}}(nt_D/\varepsilon).\]
\paragraph{Proof.} We will view the distinguisher $D$ as a deterministic algorithm that takes a string of $t_D$ bits, denoted by $r_D$, as its internal randomness. By our assumption (and dropping the absolute value sign without loss of generality), we have
\[\mathbb E_{y\sim\{0,1\}^{nk},r_D\sim\{0,1\}^{t_D}}[D(\mathrm{DP}_k(x,y);r_D)]-\mathbb E_{z\sim\{0,1\}^{nk+k},r_D\sim\{0,1\}^{t_D}}[D(z;r_D)]\geq\varepsilon.\]
Then by an averaging argument, we have
\[\Pr_{r_D}\left[\mathbb E_y[D(\mathrm{DP}_k(x,y);r_D)]-\mathbb E_z[D(z;r_D)]\geq\varepsilon/2\right]\geq\varepsilon/2.\tag{4}\]
In other words, with probability at least $\varepsilon/2$ over the internal randomness $r_D$, $D(-,r_D)$ is an $(\varepsilon/2)$-distinguisher for $\mathrm{DP}_k(x;U_{nk})$ and $U_{nk+k}$. Let us say that $r_D$ is good if this is true.
Let $\mathrm{Recon}^{(-)}:\{0,1\}^k\times\{0,1\}^r\to\{0,1\}^n$ be the reconstruction procedure in Lemma 14 that works for distinguishing parameter $\varepsilon/2$, where the randomness complexity $r$ is $\operatorname{poly}(n/\varepsilon)$. We have
\begin{align*}
&\Pr_{w\sim\{0,1\}^r,r_D\sim\{0,1\}^{t_D}}[\exists\alpha\in\{0,1\}^k\text{ such that }\mathrm{Recon}^{D(-,r_D)}(\alpha,w)=x]\\
&\geq\Pr_{w,r_D}[\exists\alpha\in\{0,1\}^k\text{ such that }\mathrm{Recon}^{D(-,r_D)}(\alpha,w)=x\mid r_D\text{ is good}]\cdot\Pr_{r_D}[r_D\text{ is good}]\\
&\geq\frac1{\operatorname{poly}(n/\varepsilon)}\cdot(\varepsilon/2).\qquad\text{(Lemma 14 and Equation (4))}
\end{align*}
The above implies that for at least $1/\operatorname{poly}(n/\varepsilon)$ fraction of the randomness $(w,r_D)$, there exists a program $M$, which takes $k$ bits for some $\alpha$ (that can depend on the randomness) and the advice of $\beta$ used by $D$, such that $M$ simulates $\mathrm{Recon}^{D(-,r_D)}$ on $(\alpha,w)$ and outputs $x$. Note that since $\mathrm{Recon}^D$ runs in time $t_D\cdot\operatorname{poly}(n/\varepsilon)$, our program $M$ can be made to run in time $T:=\widetilde O(t_D)\cdot\operatorname{poly}(n/\varepsilon)$. Therefore, we have
\[\mathrm{pK}_{1/\operatorname{poly}(n/\varepsilon)}^T(x\mid\beta)\leq k+O(\log(nt_D/\varepsilon)).\]
Then the lemma follows easily from Lemma 21.\hfill$\triangleleft$

\subsection*{3.4 Symmetry of information (original Item1 and its full proof)}
\paragraph{Lemma 26 (Symmetry of information for $\mathrm{pK}^t$ and $\mathrm{rK}^t$ under the average-case easiness of NP), Item 1.} If $\mathrm{DistNP}\subseteq\mathrm{AvgBPP}$, then there exist polynomials $p$ and $p_0$ such that for all sufficiently large $n,m\in\mathbb N$, all $t\geq p_0(n,m)$, and all $0\leq\tau\leq t$ the following holds: for every $x\in\{0,1\}^n$ and every family $\{y_r\in\{0,1\}^m\mid r\in\{0,1\}^\tau\}$, with probability at least $9/10$ over $r\in\{0,1\}^\tau$,
\[\mathrm{pK}^t(x,y_r\mid r)>\mathrm{pK}^{p(t)}(x\mid r)+\mathrm{pK}^{p(t)}(y_r\mid x,r)-\log p(t).\]
Moreover, the special case without conditioning on $r$ also holds for both items above, i.e., when $r$ is the empty string (for $\tau=0$) and $y$ is an arbitrary single string.
\paragraph{Proof of Item 1.} Define a language
\[L:=\{(u,v,w,w',1^s)\mid\exists M\in\{0,1\}^s,M(w,w')\text{ prints }uv\text{ in }|w|\text{ steps, and }s=|u|+|v|-10\}.\]
Note that $L\in\mathrm{NP}$. Define a distribution family $\mathcal D=\{D_{\langle nk+k,mk'+k',2t,\tau,s\rangle}\}$ as follows: sample $u\sim U_{nk+k}$, $v\sim U_{mk'+k'}$, $w\sim U_{2t}$, and $w'\sim U_\tau$, and then output $(u,v,w,w',1^s)$. Note that the ensemble $\mathcal D$ is in PSamp under an appropriate encoding of its defining tuple (for simplicity, we omit the padding in the output of $\mathcal D$; see Section 2.1). Using the assumption that $\mathrm{DistNP}\subseteq\mathrm{AvgBPP}$, it follows that $(L,\mathcal D)\in\mathrm{AvgBPP}$. Let $B$ be a randomized algorithm for $(L,\mathcal D)$ as in Lemma 6.
Let $x\in\{0,1\}^n$ and $\{y_r\}_{r}\subseteq\{0,1\}^m$ be given (for sufficiently large $n$ and $m$). Let $k,k'\in\mathbb N$ be arbitrary parameters such that $k\leq n$ and $k'\leq m$. Observe that there exists a polynomial $p_0$ such that for any $t\geq p_0(n,m)$, some constant $d$, and all $z\in\{0,1\}^{nk}$, $z'\in\{0,1\}^{mk'}$, and $r\in\{0,1\}^\tau$,
\[\mathrm{pK}^{2t}(\mathrm{DP}_k(x;z),\mathrm{DP}_{k'}(y_r;z')\mid r)\leq\mathrm{pK}^t(x,y_r\mid r)+|z|+|z'|+d\log t.\tag{6}\]
Here $p_0(n,m)$ reflects the time required to deterministically compute $(\mathrm{DP}_k(x;z),\mathrm{DP}_{k'}(y_r;z'))$ given $x,y_r,z,z'$, and $d\log t$ bits of information to delineate $x$ from $y_r$.
Let $t\geq p_0(n,m)$. By a counting argument, for independently random $u,v,w$, and $w'$,
\[\Pr_{u,v,w,w'}[(u,v,w,w',1^s)\in L]\leq\frac{2^s\cdot2^{|w|+|w'|}}{2^{|u|+|v|+|w|+|w'|}}=2^{s-|u|-|v|}=2^{-10},\]
where the last equality is by the definition of $L$ requiring that $s=|u|+|v|-10$. Then
\begin{align*}
\Pr_{u,v,w,w',r_B}[B(u,v,w,w',1^s)=1]
&\leq\Pr_{u,v,w,w'}[(u,v,w,w',1^s)\in L]\\
&\quad+\Pr_{u,v,w,w',r_B}[B(u,v,w,w',1^s)\neq L(u,v,w,w',1^s)]\\
&\leq2^{-10}+(3/n).\tag{7}
\end{align*}
By Markov's inequality, for at least $9/10$ fraction of random strings $r\in\{0,1\}^\tau$, we get
\[\Pr_{u,v,w,r_B}[B(u,v,w,r,1^s)=1]\leq10\cdot(2^{-10}+(3/n))\leq1/10.\tag{8}\]
Fix any such $r$ so that Equation (8) holds. This also fixes the string $y_r$.
Next we show, using a hybrid argument, that $B(-,U_{2t},r,1^s)$ cannot distinguish between the uniform distribution and the distribution $(\mathrm{DP}_k(x;z),\mathrm{DP}_{k'}(y_r;z'))$, for random independent $z,z'$, where $k\approx\mathrm{pK}^{p_{\mathrm{DP}}(t)}(x\mid r)$ and $k'\approx\mathrm{pK}^{p_{\mathrm{DP}}(t)}(y_r\mid x,r)$. This will imply that
\[\Pr_w[(\mathrm{DP}_k(x;z),\mathrm{DP}_{k'}(y_r;z'),w,r,1^s)\in L]<2/3\]
for some $z,z'$, yielding the desired lower bound on $\mathrm{pK}^{2t}(\mathrm{DP}_k(x;z),\mathrm{DP}_{k'}(y_r;z')\mid r)$ via our choice of $s$ and the definition of $L$. We give the details of the hybrid argument next.
First, toward a contradiction, suppose
\[\Pr_{z,v,w,r_B}[B(\mathrm{DP}_k(x;z),v,w,r,1^s)=1]>1/2.\]
In this case, comparing with Equation (8), we get a randomized distinguisher (with advice) for $\mathrm{DP}_k(x;U_{nk})$, defined by sampling $v\sim U_{mk'+k'}$ and $w\sim U_{2t}$, and outputting $B(-,v,w,r,1^s)$. By Lemma 22,
\[\mathrm{pK}^{q(t)}(x\mid r)\leq k+\log q(t)\tag{9}\]
for some polynomial $q$ such that $q(t)\geq p_{\mathrm{DP}}(t_B\cdot3\cdot n)$ whenever $t\geq n+m$ and $s\leq n^3$, where $p_{\mathrm{DP}}$ is the polynomial from Lemma 22 and $t_B$ denotes the time required to compute $B(-,v,w,r,1^s)$.
Define $k:=\mathrm{pK}^{q(t)}(x\mid r)-\log q(t)-1$ so that Equation (9) does not hold. Assume for now that $k>0$. Hence,
\[\Pr_{z,v,w,r_B}[B(\mathrm{DP}_k(x;z),v,w,r,1^s)=1]\leq1/2.\tag{10}\]
Again, toward a contradiction, suppose that for all $z,z'$,
\[\Pr_w[(\mathrm{DP}_k(x;z),\mathrm{DP}_{k'}(y_r;z'),w,r,1^s)\in L]\geq2/3.\]
By definition of $B$, this implies that
\[\Pr_{z,z',w,r_B}[B(\mathrm{DP}_k(x;z),\mathrm{DP}_{k'}(y_r;z'),w,r,1^s)=1]\geq(2/3)(1-1/n)>5/8.\]
In this case, comparing with Equation (10), we get a randomized distinguisher (with advice) $B'$ for $\mathrm{DP}_{k'}(y_r;U_{mk'})$, defined by sampling $z\sim U_{nk}$, $w\sim U_{2t}$, and outputting $B(\mathrm{DP}_k(x;z),-,w,r,1^s)$. By Lemma 22,
\[\mathrm{pK}^{q'(t)}(y_r\mid x,r)\leq k'+\log q'(t)\tag{11}\]
for some polynomial $q'$ with $q'(t)\geq p_{\mathrm{DP}}(t_{B'}\cdot8\cdot m)$ whenever $t\geq n+m$ and $s\leq n^3$, where $t_{B'}$ denotes the time required to compute $B'$.
We now choose $k':=\mathrm{pK}^{q'(t)}(y_r\mid x,r)-\log q'(t)-1$ so that Equation (11) does not hold. Assume for now that $k'>0$. Hence, there exist $z$ and $z'$ such that
\begin{align*}
&\Pr_w[(\mathrm{DP}_k(x;z),\mathrm{DP}_{k'}(y_r;z'),w,r,1^s)\in L]\\
&=\Pr_w[\exists M\in\{0,1\}^s,M(w,r)\text{ outputs }(\mathrm{DP}_k(x;z),\mathrm{DP}_{k'}(y_r;z'))\text{ within }|w|\text{ steps}]<2/3,
\end{align*}
which implies that
\[\mathrm{pK}^{2t}(\mathrm{DP}_k(x;z),\mathrm{DP}_{k'}(y_r;z')\mid r)>s.\]
For this choice of $z,z'$, by definition of $s$ we get that $s=|u|+|v|-10=|z|+k+|z'|+k'-10$, and so
\[\mathrm{pK}^{2t}(\mathrm{DP}_k(x;z),\mathrm{DP}_{k'}(y_r;z')\mid r)>|z|+k+|z'|+k'-10.\]
Combining this inequality with Equation (6), we get
\begin{align*}
\mathrm{pK}^t(x,y_r\mid r)
&\geq\mathrm{pK}^{2t}(\mathrm{DP}_k(x;z),\mathrm{DP}_{k'}(y_r;z')\mid r)-|z|-|z'|-d\log t\\
&>k+k'-d\log t-10.\tag{12}
\end{align*}
The definitions of $k$ and $k'$ then imply that
\begin{multline*}
\mathrm{pK}^t(x,y_r\mid r)>\mathrm{pK}^{q(t)}(x\mid r)+\mathrm{pK}^{q'(t)}(y_r\mid x,r)\\
-\log q(t)-\log q'(t)-d\log t-12.
\end{multline*}
For the polynomial $p(t):=(t^{d+1})\cdot q(t)\cdot q'(t)$ and for every $t\geq p_0(|x|,|y|)$, we get
\[\mathrm{pK}^t(x,y_r\mid r)>\mathrm{pK}^{p(t)}(x\mid r)+\mathrm{pK}^{p(t)}(y_r\mid x,r)-\log p(t),\]
as desired.
Finally, consider the case that $k\leq0$ or $k'\leq0$. If $k\leq0$, then $\mathrm{pK}^{q(t)}(x\mid r)\leq\log q(t)+1$, implying that $\mathrm{pK}^{p(t)}(x\mid r)<\log p(t)$. But then the lemma simply follows from the fact that $\mathrm{pK}^t(x,y_r\mid r)\geq\mathrm{pK}^{p(t)}(y_r\mid x,r)$. Similarly, if $k'\leq0$, then $\mathrm{pK}^{p(t)}(y_r\mid x,r)<\log p(t)$, and the lemma follows from the fact that $\mathrm{pK}^t(x,y_r\mid r)\geq\mathrm{pK}^{p(t)}(x\mid r)$.\hfill$\triangleleft$
\subsection*{3.7 Probabilistic computational depth (original Item1 and its full proof)}
The next lemma is a straightforward adaptation of an argument from [11]. We provide a proof for completeness.
\paragraph{Lemma 31 (Computational Depth Upper Bound), Item 1.} For all $x\in\{0,1\}^n$ and $r\in\{0,1\}^*$, and for all non-decreasing polynomials $p_0$ and $p$, if $n$ is large enough then there exists a time bound $t_1$ such that $p_0(n)\leq t_1\leq2^{n/\log n}$ and
\[\mathrm{pK}^{t_1}(x\mid r)-\mathrm{pK}^{p(t_1)}(x\mid r)\leq O(n/\log n).\]
\paragraph{Proof.} We will start by proving Item 1. Given $x\in\{0,1\}^n$ and polynomials $p_0$ and $p$, define the polynomial $\tau:=p\circ p_0$. For an integer $I\geq1$, consider the following telescoping sum:
\begin{align*}
\mathrm{pK}^{\tau(n)}(x\mid r)-\mathrm{pK}^{\tau^{I+1}(n)}(x\mid r)
&=(\mathrm{pK}^{\tau(n)}(x\mid r)-\mathrm{pK}^{\tau^2(n)}(x\mid r))\\
&\quad+(\mathrm{pK}^{\tau^2(n)}(x\mid r)-\mathrm{pK}^{\tau^3(n)}(x\mid r))\\
&\quad+\cdots+(\mathrm{pK}^{\tau^I(n)}(x\mid r)-\mathrm{pK}^{\tau^{I+1}(n)}(x\mid r)),
\end{align*}
where $\tau^i(-)$ denotes the composition of $\tau$ with itself $i$ times. For any choice of $x,p_0$, and $p$ as in the statement of the lemma, $\mathrm{pK}^{\tau(n)}(x\mid r)\leq n+d$, for some universal constant $d\geq0$; hence, the above sum is at most $n+d$. By averaging, there is some index $i_0\in[I]$ such that
\[\mathrm{pK}^{\tau^{i_0}(n)}(x\mid r)-\mathrm{pK}^{\tau^{i_0+1}(n)}(x\mid r)\leq(n+d)/I.\]
For this $i_0$, define $t_1:=\tau^{i_0}(n)$. Note that $t_1\geq\tau(n)\geq p_0(n)$, since $i_0\geq1$ and $p(\ell)\geq\ell$ for every input $\ell$. Letting $c\in\mathbb N$ be such that $\tau(n)\leq n^c$ for sufficiently large $n$, define $I:=\log_c(n/(\log_2n)^2)$. Then $t_1\leq n^{c^I}=2^{n/\log n}$. Moreover,
\begin{align*}
\mathrm{pK}^{t_1}(x\mid r)-\mathrm{pK}^{p(t_1)}(x\mid r)&\leq\mathrm{pK}^{t_1}(x\mid r)-\mathrm{pK}^{\tau(t_1)}(x\mid r)\\
&\leq O(n/\log n),
\end{align*}
as desired.

% Faithful official CCC2022 article16 original pKt cases of Lemmas33–34 and Theorem36.
\section*{4 Probabilistic Worst-Case to Average-Case Reductions (AM case)}
\subsection*{4.1 Auxiliary lemmas}
Given a language $L\in\mathrm{NP}$ with a corresponding verifier $V$, and $x\in L$, let $y_x$ be the lexicographically first witness that $x\in L$. It is not hard to compute $y_x$ from $x$ if an NP oracle is available, by performing a standard search-to-decision reduction. For this reason, the time-bounded Kolmogorov complexity of the pair $(x,y_x)$ is essentially that of $x$, in the presence of an NP oracle. We will show, more generally, that the oracle can be eliminated under an appropriate average-case easiness assumption, even for the following randomized versions of $\Sigma_\ell^\mathrm P$, denoted $\mathrm{BP}_\delta\circ\Sigma_\ell^\mathrm P$.
\paragraph{Definition 32 ($\mathrm{BP}_\delta\circ\Sigma_\ell^\mathrm P$).} A language $L$ is in the class $\mathrm{BP}_\delta\circ\Sigma_\ell^\mathrm P$ if there is a $\mathrm{BP}\circ\Sigma_\ell\circ\mathrm P$ formula
\[\varphi(x)=\mathrm{BP}_r\exists y_1\forall y_2\ldots Qy_\ell\ R(x,y_1,y_2,\ldots,y_\ell,r)\]
for a polytime predicate $R$ and all string variables $y_1,\ldots,y_\ell,r$ of length $\operatorname{poly}(|x|)$, such that, for all $x\in\{0,1\}^n$,
\begin{align*}
x\in L&\Longrightarrow\Pr_r[\exists y_1\forall y_2\ldots Qy_\ell\ R(x,y_1,y_2,\ldots,y_\ell,r)]\geq\delta,\text{ and}\\
x\notin L&\Longrightarrow\Pr_r[\forall y_1\exists y_2\ldots\overline Qy_\ell\ \neg R(x,y_1,y_2,\ldots,y_\ell,r)]\geq\delta.
\end{align*}
\paragraph{Lemma 33 (Oracle Elimination), original Item 1.} For any $\ell\geq1$, suppose $\mathrm{Dist}\Sigma_{\ell+1}^{\mathrm P}\subseteq\mathrm{AvgBPP}$. For an arbitrary $L\in\mathrm{BP}_\delta\circ\Sigma_\ell^{\mathrm P}$, let $x\in L_n$ be sufficiently large, let $r\in\{0,1\}^{n^c}$ be any random string under the BP quantifier such that there exists a witness for the left-most $\exists$ quantifier in the definition of $L$, and let $y_{x,r}\in\{0,1\}^{n^c}$ be the lexicographically first such $L$-witness for $x$ and randomness $r$, for some constant $c>0$.
There exist polynomials $q$ and $q_0$ such that, for all sufficiently large $n\in\mathbb N$, all $x\in L_n$, and all $t\geq q_0(n)$, we have with probability at least $\delta-1/10$ over $r\in\{0,1\}^{n^c}$ that an $L$-witness $y_{x,r}$ exists, and
\[\mathrm{pK}^{q(t)}(x,y_{x,r}\mid r)\leq\mathrm{pK}^t(x\mid r)+\log q(t).\]
Moreover, for the case of no BP quantifier, i.e., for $L\in\Sigma_\ell^{\mathrm P}$, we get the same conclusions (with probability $1$) without conditioning on $r$.
\paragraph{Proof of Item 1.} Define
\begin{multline*}
L':=\{(\mathrm{DP}_k(x,y;z),w,w',1^k,1^s,1^n)\mid\exists M\in\{0,1\}^s,\\
M^{\Sigma_\ell^{\mathrm P}}(w,w')\text{ prints }(x,y)\in\{0,1\}^{n+n^c}\text{ within }|w|\text{ steps, where }|w'|=n^c,\text{ and }s=k-10\}.
\end{multline*}
More formally, in this and subsequent proofs, the above notation should be interpreted as follows: a string $(u,w,w',1^k,1^s,1^n)$ belongs to $L'$ iff there exist some $(x,y)\in\{0,1\}^{n+n^c}$ and $M\in\{0,1\}^s$ such that $\mathrm{DP}_k(x,y;z)=u$, where $u=(z,\alpha)$ for some $z\in\{0,1\}^{(n+n^c)k}$ and $\alpha\in\{0,1\}^k$, and $M^{\Sigma_\ell^{\mathrm P}}(w,w')$ outputs $(x,y)$ within $|w|$ steps, for $|w'|=n^c$.
Note that $L'\in\mathrm{NP}^{\Sigma_\ell^{\mathrm P}}$. Define a distribution family $\mathcal D=\{D_{\langle(n+n^c)k+k,2t,n^c,k,s,n\rangle}\}$ as follows: sample $u\sim U_{(n+n^c)k+k}$, $w\sim U_{2t}$, $w'\sim U_{n^c}$, and output $(u,w,w',1^k,1^s,1^n)$, for $s=k-10$. By assumption, $(L',\mathcal D)\in\mathrm{AvgBPP}$. Let $B$ be a randomized algorithm for $(L',\mathcal D)$ as in Lemma 6.
For all (sufficiently large) $n,k$, and $s$, and for independent uniformly random $u,w$, and $w'$,
\begin{align*}
&\Pr_{u,w,w'}[(u,w,w',1^k,1^s,1^n)\in L']\\
&=\Pr_{u,w,w'}[\exists(x,y)\in\{0,1\}^{n+n^c},\exists z\in\{0,1\}^{(n+n^c)k},\exists M\in\{0,1\}^s,\\
&\qquad M^{\Sigma_\ell^{\mathrm P}}(w,w')=(x,y)\wedge u=\mathrm{DP}_k(x,y;z)]\\
&\leq\frac{2^{s+|w|+|w'|+|z|}}{2^{k+|w|+|w'|+|z|}}=2^{s-k}=2^{-10},
\end{align*}
where the inequality is by a union bound and a counting argument, and the last equality by the condition that $s=k-10$ in the definition of $L'$. Hence,
\begin{align*}
\Pr_{u,w,w',r_B}[B(u,w,w',1^k,1^s,1^n)=1]
&\leq\Pr_{u,w,w'}[(u,w,w',1^k,1^s,1^n)\in L']\\
&\qquad+\Pr_{u,w,w',r_B}[B(u,w,w',1^k,1^s,1^n)\neq L'(u,w,w',1^k,1^s,1^n)]\\
&\leq2^{-10}+(3/n).
\end{align*}
By Markov's inequality, for at least $9/10$ fraction of strings $r$,
\[\Pr_{u,w,r_B}[B(u,w,r,1^k,1^s,1^n)=1]\leq10\cdot(2^{-10}+(3/n))\leq1/10.\tag{18}\]
Let $x\in L_n$ be arbitrary. By definition of the $\mathrm{BP}_\delta$ quantifier, for at least $\delta$ fraction of random strings $r\in\{0,1\}^{n^c}$, there exists an $L$-witness, and hence, the lexicographically first $L$-witness $y_{x,r}$. By the above, for at least $\delta-1/10$ of random $r$, we have that both a witness $y_{x,r}$ exists and Equation (18) holds. Fix any such $r$.
Observe that for some polynomial $q_0$ (dependent on $L$) and some constant $d$, for any $t\geq q_0(n)$,
\[\mathrm{pK}^{2t,\Sigma_\ell^{\mathrm P}}(x,y_{x,r}\mid r)\leq\mathrm{pK}^t(x\mid r)+d\log n=:s(x,r).\tag{19}\]
In particular, $q_0(n)$ reflects the time required to deterministically compute $y_{x,r}$, given $x$ and $r$, by a search-to-decision procedure for $L$ using a $\Sigma_\ell^{\mathrm P}$-oracle.
Define $s:=s(x,r)$ (which determines $k=s+10$), and let $t\geq q_0(n)$. By the definitions of $\mathrm{pK}^{2t,\Sigma_\ell^{\mathrm P}}$ and $L'$, for every $z\in\{0,1\}^{(n+n^c)k}$,
\[\Pr_w[(\mathrm{DP}_k(x,y_{x,r};z),w,r,1^k,1^s,1^n)\in L']\geq2/3.\]
Then by the definition of $B$,
\begin{align*}
&\Pr_{z,w,r_B}[B(\mathrm{DP}_k(x,y_{x,r};z),w,r,1^k,1^s,1^n)=1]\\
&\geq\Pr_{z,w}[(\mathrm{DP}_k(x,y_{x,r};z),w,r,1^k,1^s,1^n)\in L']\cdot\Pr_{r_B}[B(\omega)=L'(\omega)\mid\omega\in L']\\
&\geq\frac23(1-(1/n))\geq2/3-o(1).\tag{20}
\end{align*}
Comparing Equation (18) with Equation (20), it is clear that $B(-,U_{2t},r,1^k,1^s,1^n)$ is a randomized algorithm (with advice) that $(1/2)$-distinguishes $\mathrm{DP}_k(x,y_{x,r};U_{(n+n^c)k})$ from uniform. Lemma 22 implies that
\begin{align*}
\mathrm{pK}^{p'(t)}(x,y_{x,r}\mid r,s,n,t)&\leq k+\log p'(t)\tag{21}\\
&=s+10+\log p'(t)\\
&=\mathrm{pK}^t(x\mid r)+d\log n+\log p'(t)+10,
\end{align*}
for some polynomial $p'$ with $p'(t)\geq p_{\mathrm{DP}}(t_B\cdot3\cdot(n+n^c))$, where $p_{\mathrm{DP}}$ is the polynomial from Lemma 22 and $t_B$ denotes the time required to compute $B(-,U_{2t},r,1^k,1^s,1^n)$. As the advice $(s,n,t)$ can be encoded with $\log s+\log n+\log t+O(\log\log n)\leq(3.1)\cdot\log t$ bits, it follows that $\mathrm{pK}^{q(t)}(x,y_{x,r}\mid r)\leq\mathrm{pK}^t(x\mid r)+\log q(t)$ for the polynomial $q(t):=t^{d+4}\cdot p'(t)$.
The ``Moreover'' statement follows the same argument, dropping any mention of $w'$ (and $r$).\hfill$\triangleleft$
\paragraph{Lemma 34 (Witness Compression), original Item 1.} For any $\ell\geq1$, suppose $\mathrm{Dist}\Sigma_{\ell+1}^{\mathrm P}\subseteq\mathrm{AvgBPP}$. For an arbitrary $L\in\mathrm{BP}_\delta\circ\Sigma_\ell^{\mathrm P}$, let $x\in L_n$ be sufficiently large, $r\in\{0,1\}^{n^c}$ be any good random string, and let $y_{x,r}\in\{0,1\}^{n^c}$ be the lexicographically first $L$-witness for $x$ with respect to randomness $r$, for some constant $c>0$.
Let $p_0,p$ and $q_0,q$ be the polynomials from Items 1 of Lemmas 26 and 33 respectively. Then for every $t\geq\max\{p_0(n),q_0(n+n^c)\}$, with probability at least $\delta-1/5$ over $r$, an $L$-witness $y_{x,r}$ exists, and
\[\mathrm{pK}^{p(q(t))}(y_{x,r}\mid x,r)\leq(\mathrm{pK}^t(x\mid r)-\mathrm{pK}^{p(q(t))}(x\mid r))+2\log p(q(t)).\]
Moreover, for $L\in\Sigma_\ell^{\mathrm P}$, we get the same conclusions without conditioning on $r$.
\paragraph{Proof.} We start by proving Item 1. By Item 1 of Lemma 33 and Lemma 26, we get by a union bound that, with probability at least $\delta-2/10$ over $r$,
\[\mathrm{pK}^{q(t)}(x,y_{x,r}\mid r)\leq\mathrm{pK}^t(x\mid r)+\log q(t),\]
and
\[\mathrm{pK}^{q(t)}(x,y_{x,r}\mid r)>\mathrm{pK}^{p(q(t))}(x\mid r)+\mathrm{pK}^{p(q(t))}(y_{x,r}\mid x,r)-\log p(q(t)).\]
Combining the previous two inequalities,
\begin{align*}
\mathrm{pK}^{p(q(t))}(y_{x,r}\mid x,r)
&<\mathrm{pK}^{q(t)}(x,y_{x,r}\mid r)-\mathrm{pK}^{p(q(t))}(x\mid r)+\log p(q(t))\\
&\leq(\mathrm{pK}^t(x\mid r)-\mathrm{pK}^{p(q(t))}(x\mid r))+2\log p(q(t)).
\end{align*}
The ``Moreover'' part follows by applying the ``no random $r$'' versions of Lemma 33 and Lemma 26.
\subsection*{4.2 Case of AM}
\paragraph{Theorem 36.} If $\mathrm{Dist}\Sigma_2^{\mathrm P}\subseteq\mathrm{AvgBPP}$, then $\mathrm{AM}\subseteq\mathrm{BPTIME}[2^{O(n/\log n)}]$.
\paragraph{Proof.} Consider an arbitrary $L\in\mathrm{AM}$ given by a $\mathrm{BP}\circ\Sigma_1^{\mathrm P}$ formula such that, for every $x\in\{0,1\}^n$,
\begin{align*}
x\in L&\Longrightarrow\mathrm{BP}_{(1-2^{-n})}r\ \exists y\ V(x,y,r),\\
x\notin L&\Longrightarrow\mathrm{BP}_{(1-2^{-n})}r\ \forall y\ \neg V(x,y,r),
\end{align*}
for a polytime verifier $V(x,y,r)$, where $|r|,|y|\leq n^c$ for some constant $c>0$, and the notation $\mathrm{BP}_{(\delta)}r\ \varphi(-,r)$, for a formula $\varphi$, means $\Pr_r[\varphi(-,r)]\geq\delta$.
For a given (sufficiently large) $x\in L_n$, and a good random string $r$, let $y_{x,r}$ be the lexicographically first $L$-witness for $x,r$. By Lemma 34, for some constant $a>0$, and for all large enough time bounds $t$, we have, with probability at least $1-2^{-n}-1/5$ over random strings $r$, that
\[\mathrm{pK}^{t^a}(y_{x,r}\mid x,r)\leq(\mathrm{pK}^t(x\mid r)-\mathrm{pK}^{t^a}(x\mid r))+(2a)\log t.\]
By Lemma 31, there exists a time bound $t\leq2^{O(n/\log n)}$ such that $\mathrm{pK}^t(x\mid r)-\mathrm{pK}^{t^a}(x\mid r)\leq bn/\log n$, for some universal constant $b>0$. Hence, for some universal constant $d>0$ (independent of $x,r$),
\[\mathrm{pK}^{2^{dn/\log n}}(y_{x,r}\mid x,r)\leq dn/\log n.\]
A randomized algorithm for $L$ is now obvious:
\begin{itemize}
\item Given $x\in\{0,1\}^n$, sample uniformly random strings $r\in\{0,1\}^{|x|^c}$ and $w\in\{0,1\}^{2^{dn/\log n}}$, and then exhaustively search over all $dn/\log n$-bit machines $M$, running each $M(w,x,r)$ for $2^{dn/\log n}$ steps to produce a candidate witness $y$, accepting if $V(x,y,r)$ accepts.
\end{itemize}
The runtime $2^{O(n/\log n)}$ is clear. For correctness, for every $x\in L_n$, we know there is at least a $(4/5-2^{-n})$ fraction of good random strings $r$ for which a witness $y_{x,r}$ will be found by our algorithm, with probability at least $2/3$ over random $w$'s (by the definition of pK). So our algorithm accepts $x\in L$ with probability at least $(2/3)(4/5)-o(1)>1/2$. On the other hand, for any $x\notin L$, our algorithm may find a witness $y$ and accept $x$ for at most $2^{-n}$ fraction of $r$'s.\hfill$\triangleleft$

% Faithful official CCC2022 article16 additional pKt coding/approximation and agnostic-learning definitions.
\section*{2.6 Agnostic learning model (author definitions)}
Consider the problem of learning an unknown concept $f:X\to\{0,1\}$ over a finite domain $X$. We will always let $X=\{0,1\}^n$, for some $n\geq1$. A concept class $\mathcal C$ is a collection of concepts of this form. For a fixed $n$, we let $\mathcal C_n$ denote $\mathcal C\cap\{f:\{0,1\}^n\to\{0,1\}\}$. For a size function $s:\mathbb N\to\mathbb N$, we let
\[\mathrm{SIZE}[s]=\{f:\{0,1\}^n\to\{0,1\}\mid n\in\mathbb N\text{ and }f\text{ admits a size-}s(n)\text{ circuit}\}\]
denote the concept class of Boolean circuits of size (number of gates) at most $s(n)$ over $n$ input variables. For definiteness, we consider circuits over AND and OR gates of fan-in $2$ and NOT gates.
A randomized Boolean function $f$ maps an input $x$ to a distribution $D$ supported over $\{0,1\}$. Given a distribution $D_n$ supported over $\{0,1\}^n$, a function $h:\{0,1\}^n\to\{0,1\}$, and a (possibly randomized) Boolean function $f:\{0,1\}^n\to\{0,1\}$, we let
\[\operatorname{err}_{D_n}(h,f)=\Pr_{f,x\sim D_n}[h(x)\neq f(x)].\]
We will consider the (agnostic) learnability of a concept class $\mathcal C$ with respect to a class $\mathcal D$ of ensembles $D=\{D_n\}_{n\geq1}$ of distributions $D_n$. We use $\mathcal D_n$ to denote $\{D_n\mid D\in\mathcal D\}$. For simplicity, we refer to $\mathcal D$ as a class of distributions.
Given a randomized function $f:\{0,1\}^n\to\{0,1\}$ and a distribution $D_n$, the learning algorithm has access to an example oracle $\mathrm{EX}(f,D_n)$ that behaves as follows: each query to $\mathrm{EX}(f,D_n)$ returns an independent and identically distributed pair $(x,b)$, where $x$ is sampled according to $D_n$ and $b=f(x)$. The sample complexity of the learning algorithm is the number of queries made to $\mathrm{EX}(f,D_n)$.
Before formalising the agnostic learning model, we provide an informal overview. The learning algorithm knows the class $\mathcal D$ of distributions and the concept class $\mathcal C$. It is given access to $\mathrm{EX}(f,D_n)$ for some unknown $D\in\mathcal D_n$ and an arbitrary (possibly randomized) function $f:\{0,1\}^n\to\{0,1\}$. The goal of the learning algorithm is to produce, with high probability over its internal randomness and queries to $\mathrm{EX}(f,D)$, a hypothesis $h$ such that $\operatorname{err}_D(h,f)\leq\operatorname{opt}_{\mathcal C_n,D,f}+\varepsilon$, where
\[\operatorname{opt}_{\mathcal C_n,D,f}=\min_{c\in\mathcal C_n}\operatorname{err}_D(c,f).\]
In other words, the learning algorithm is required to output an efficient representation of a hypothesis $h:\{0,1\}^n\to\{0,1\}$ that is almost as accurate as the best concept in $\mathcal C_n$ with respect to the pair $(D,f)$. (Note that $h$ is not required to be a function in $\mathcal C$.)
\paragraph{Definition 15 (Agnostic PAC learning [20]).} Let $\mathcal C$ be a concept class, and let $\mathcal D$ be a class of distributions. We say that a randomized algorithm $A$ agnostically learns $\mathcal C$ on $\mathcal D$ if the following holds. For every $n\geq1$, distribution $D\in\mathcal D_n$, and randomized function $f:\{0,1\}^n\to\{0,1\}$, when $A$ is given oracle access to $\mathrm{EX}(f,D)$ and receives as input $n,\varepsilon>0$, and $\delta>0$,
\[\Pr_{A,\mathrm{EX}(f,D)}[A^{\mathrm{EX}(f,D)}\text{ outputs a hypothesis }h\text{ such that }\operatorname{err}_D(h,f)\leq\operatorname{opt}_{\mathcal C_n,D,f}+\varepsilon]\geq1-\delta.\]
We remark that an equivalent way of formulating the agnostic learning model is by considering distributions $D$ supported over $\{0,1\}^{n+1}$. In this case, there is no need to refer to randomized functions, and one measures the error of a hypothesis $h$ via $\operatorname{err}_D(h)=\Pr_{(x,b)\sim D}[h(x)\neq b]$, where $x\in\{0,1\}^n$ and $b\in\{0,1\}$.
\section*{Additional author pKt properties}
\paragraph{Lemma 18.} There is a universal constant $c>0$ such that the following holds. For every time bound $t\in\mathbb N$ and $x\in\{0,1\}^n$,
\[\mathrm K(x\mid t)\leq\mathrm{pK}^t(x)+c\log n.\]
\paragraph{Proof.} Let $\ell:=\mathrm{pK}^t(x)$. Since $x$ is of length $n$, we can represent $\ell$ using $O(\log n)$ bits. Given $t,\ell$ and $n$, it is possible to sample $x$ with probability $\mu\geq\frac23\cdot2^{-\ell}$ by randomly guessing $w\sim\{0,1\}^t$ and $M\sim\{0,1\}^\ell$ then simulating $M(w)$ for $t$ steps. Consequently, by the coding theorem for (time-unbounded) Kolmogorov complexity [24], we get that $\mathrm K(x\mid t)\leq\log(1/\mu)+O(\log n)\leq\mathrm{pK}^t(x)+O(\log n)$.\hfill$\triangleleft$
\paragraph{Lemma 20 (Probabilistic Incompressibility).} For any string $y\in\{0,1\}^*$, time bound $t\in\mathbb N$ (including $t=\infty$), $\delta\in(0,1]$, and positive integer $\alpha$, we have
\[\Pr_{x\sim\{0,1\}^n}(\mathrm{pK}_\delta^t(x\mid y)\leq n-\alpha)\leq\frac1{2^{\alpha-1}\cdot\delta}.\]
\paragraph{Proof.} For the sake of contradiction, suppose the statement of the lemma does not hold. Then by the definition of $\mathrm{pK}_\delta^t$, we have
\begin{multline*}
\mathbb E_{x\sim\{0,1\}^n,w\sim\{0,1\}^t}[\exists M_{(x,w)}\in\{0,1\}^{\leq n-\alpha}\text{ such that }M_{(x,w)}(w,y)\text{ outputs }x\text{ within }t\text{ steps}]\\
>1/2^{\alpha-1}.
\end{multline*}
By averaging, there is a way to fix $w\in\{0,1\}^t$ that at least preserves this expectation, yielding
\begin{multline*}
\Pr_{x\sim\{0,1\}^n}[\exists M_{(x,w)}\in\{0,1\}^{\leq n-\alpha}\text{ such that }M_{(x,w)}(w,y)\text{ outputs }x\text{ within }t\text{ steps}]\\
>1/2^{\alpha-1}.
\end{multline*}
However, by counting, the above probability is at most $2^{n-\alpha+1}/2^n=1/2^{\alpha-1}$. A contradiction.\hfill$\triangleleft$
\paragraph{Lemma 24 (Unconditional Source Coding for $\mathrm{pK}^t$ [34]).} There exists a polynomial $p$ such that for any $T,a:\mathbb N\to\mathbb N$, $n\in\mathbb N$, $D_n\in\mathrm{Samp}[T(n)]/a(n)$, and $x\in\operatorname{Supp}(D_n)$,
\[\mathrm{pK}^{p(T(n))}(x)\leq\log(1/D_n(x))+O(\log(T(n)))+a(n).\]
For a distribution $D_n$ supported over $X$ and an integer $m\geq1$, we let $D_n^m:=D_n\times\ldots\times D_n$ be the probability distribution supported over $X^m$ obtained by taking the product of $m$ independent copies of $D_n$.
As a consequence of Lemma 24, we obtain the following result.
\paragraph{Lemma 25 (Source Coding for $m$ copies of $D_n$).} There exists a polynomial $p$ such that for any $T,a:\mathbb N\to\mathbb N$, $n,m\in\mathbb N$, $D_n\in\mathrm{Samp}[T(n)]/a(n)$, and $x\in\operatorname{Supp}(D_n^m)$,
\[\mathrm{pK}^{p(T(n),m)}(x)\leq\log(1/D_n^m(x))+O(\log(T(n)))+a(n)+O(\log(m)).\]
\paragraph{Proof.} Since $D_n$ can be sampled in time $T(n)$ using $a(n)$ bits of advice, it is possible to sample from $D_n^m$ in time $\operatorname{poly}(T(n),m)$ using $a(n)+O(\log m)$ bits of advice, where the extra advice is used to encode $m$. The result now immediately follows from Lemma 24.\hfill$\triangleleft$
\subsection*{3.5 Approximating $\mathrm{pK}^t$ under randomized average-case easiness}
\paragraph{Lemma 27.} If $\mathrm{DistNP}\subseteq\mathrm{AvgBPP}$, then there exists a polynomial $\tau$ such that the following promise problem is in promiseBPP:
\begin{align*}
\Pi_{\mathrm{YES}}&:=\{(x,1^s,1^t)\mid\mathrm{pK}^t(x)\leq s\},\\
\Pi_{\mathrm{NO}}&:=\{(x,1^s,1^t)\mid\mathrm{pK}^{\tau(t)}(x)>s+\log\tau(t)\}.
\end{align*}
\paragraph{Proof.} Let $k=s+\log n$. Consider the language
\[L:=\{(\mathrm{DP}_k(x;z),w,1^s,1^n)\mid\exists M\in\{0,1\}^s,M(w)\text{ outputs }x\in\{0,1\}^n\text{ within }|w|\text{ steps}\}.\]
Note that $L\in\mathrm{NP}$. Define a distribution family $\mathcal D=\{D_{\langle nk+k,t,s,n\rangle}\}$, each member of which does the following: sample $u\sim U_{nk+k}$ and $w\sim U_t$, and then output $(u,w,1^s,1^n)$. By assumption, $(L,\mathcal D)\in\mathrm{AvgBPP}$. Let $B$ be a randomized heuristic algorithm for $(L,\mathcal D)$ as described in Lemma 6.
Now, define an algorithm $B'$:
\begin{itemize}
\item On input $(x,1^s,1^t)$ with $x\in\{0,1\}^n$, sample $z\sim U_{nk}$ and $w\sim U_t$, and then output $B(\mathrm{DP}_k(x;z),w,1^s,1^n)$.
\end{itemize}
Below, we show that $B'$ solves $(\Pi_{\mathrm{YES}},\Pi_{\mathrm{NO}})$ correctly with high probability in the worst case.
First, consider the case that $(x,1^s,1^t)\in\Pi_{\mathrm{YES}}$. By the definitions of $L$ and pK, for any choice of $z\in\{0,1\}^{nk}$,
\[\Pr_w[(\mathrm{DP}_k(x;z),w,1^s,1^n)\in L]\geq2/3.\]
The definition of $B$ then implies that
\[\Pr_{w,z,r_B}[B(\mathrm{DP}_k(x;z),w,1^s,1^n)=1]>1/2,\]
and so
\[\Pr_{r_{B'}}(B'(x,1^s,1^t)=1)>1/2,\tag{13}\]
where $r_B$ denotes the internal randomness of $B$, and $r_{B'}=(w,z,r_B)$ that of $B'$.
Now consider the case that $(x,1^s,1^t)\in\Pi_{\mathrm{NO}}$. For a contradiction, suppose that
\[\Pr_{w,z,r_B}[B(\mathrm{DP}_k(x;z),w,1^s,1^n)=1]>1/3.\tag{14}\]
By a counting argument, for randomly selected $u$ and $w$,
\[\Pr_{u,w}[(u,w,1^s,1^n)\in L]\leq\frac{2^s\cdot2^{nk}\cdot2^{|w|}}{2^{nk+k+|w|}}=1/n,\]
where the last line follows from the definition of $k=s+\log n$. Then by definition of $B$,
\[\Pr_{u,w,r_B}[B(u,w,1^s,1^n)=1]\leq4/n=o(1).\tag{15}\]
Comparing Equations (14) and (15), it is clear that $B(-,U_t,1^s,1^n)$ $(1/4)$-distinguishes $\mathrm{DP}_k(x;U_{nk})$ from $U_{nk+k}$. Lemma 22 implies that
\[\mathrm{pK}^{p'(t)}(x)\leq k+\log p'(t)=s+\log n+\log p'(t),\]
for some polynomial $p'$ with $p'(t)\geq p_{\mathrm{DP}}(t_B\cdot4\cdot n)$, where $p_{\mathrm{DP}}$ is the polynomial from Lemma 22 and $t_B$ denotes the time required to compute $B(-,U_t,1^s,1^n)$. For the polynomial $\tau(t)=t\cdot p'(t)$, this means that $(x,1^s,1^t)$ is not in $\Pi_{\mathrm{NO}}$, which gives the desired contradiction. By definition of $B'$, we have that
\[\Pr_{r_{B'}}(B'(x,1^s,1^t)=1)\leq1/3.\tag{16}\]
By Equations (13) and (16), $B'$ yields a promiseBPP algorithm for $(\Pi_{\mathrm{YES}},\Pi_{\mathrm{NO}})$ via standard success amplification techniques.\hfill$\triangleleft$
\paragraph{Lemma 28 (Approximating $\mathrm{pK}^t$).} If $\mathrm{DistNP}\subseteq\mathrm{AvgBPP}$, then there exist a polynomial $\tau$, a constant $C\geq1$, and a randomized algorithm $\mathrm{Approx}_{\tau}\text{-}\mathrm{pK}$ that, on input $(x,1^t)$ where $x\in\{0,1\}^n$ and $t\geq Cn$, runs in time $\operatorname{poly}(n,t)$ and with probability at least $1-o(1)$ outputs an integer $\widetilde s$ such that
\[\mathrm{pK}^{\tau(t)}(x)-\log\tau(t)\leq\widetilde s\leq\mathrm{pK}^t(x).\]
\paragraph{Proof.} Consider a polynomial-time randomized algorithm $A$ that solves the promise problem from Lemma 27. Assume without loss of generality that on the inputs satisfying the promise its error is at most $1/n^2$, where $n=|x|$. Algorithm $\mathrm{Approx}_\tau\text{-}\mathrm{pK}$ runs $A$ on $(x,1^s,1^t)$ for $s=1,2,\ldots,n+\log n$, and outputs the first $\widetilde s$ such that $A(x,1^{\widetilde s},1^t)=1$.
The correctness of $\mathrm{Approx}_\tau\text{-}\mathrm{pK}$ follows by a union bound. Indeed, if $s<\mathrm{pK}^{\tau(t)}(x)-\log\tau(t)$, i.e., $\mathrm{pK}^{\tau(t)}(x)>s+\log\tau(t)$, using the promise we get that $\Pr_A[A(x,1^s,1^t)=1]\leq1/n^2$. On the other hand, if $s=\mathrm{pK}^t(x)$, which implies that $\mathrm{pK}^t(x)\leq s$ and the promise is satisfied, we have $\Pr_A[A(x,1^s,1^t)=1)]\geq1-1/n^2$. Since $\mathrm{pK}^t(x)\leq n+\log n$ if $t\geq Cn$, where $C$ is a sufficiently large constant, with high probability over the internal randomness of $\mathrm{Approx}_\tau\text{-}\mathrm{pK}$, it outputs a value $\widetilde s$ such that $\mathrm{pK}^{\tau(t)}(x)-\log\tau(t)\leq\widetilde s\leq\mathrm{pK}^t(x)$.\hfill$\triangleleft$

% Faithful official CCC2022 article16 Section4.5 and AppendixB fine-grained main results.
\section*{4.5 Fine-grained case}
Recall the quasilinear-time complexity classes $\mathrm{QL}=\mathrm{DTIME}[\widetilde O(n)]$, $\mathrm{NQL}=\mathrm{NTIME}[\widetilde O(n)]$, and the quasilinear-time analog $\mathrm{QLH}=\bigcup_{\ell\geq0}\Sigma_\ell^{\mathrm{QL}}$ of the polynomial-time hierarchy PH; see, e.g., [36] for more details. Define $\mathrm{BPQL}=\mathrm{BPTIME}[\widetilde O(n)]$, the quasilinear-time version of BPP. Also, define $\mathrm{QLSamp}$ to be the class of distribution families that are quasilinear-time samplable. We show better probabilistic worst-case to average-case reductions under fine-grained average-case easiness assumptions.
\paragraph{Theorem 45.} $\Sigma_2^{\mathrm{QL}}\times\mathrm{QLSamp}\subseteq\mathrm{AvgBPQL}\Longrightarrow\mathrm{AMTIME}[2^{O(\sqrt{n\log n})}]=\mathrm{BPTIME}[2^{O(\sqrt{n\log n})}]$.
It suffices to prove the inclusion $\mathrm{AMTIME}[2^{O(\sqrt{n\log n})}]\subseteq\mathrm{BPTIME}[2^{O(\sqrt{n\log n})}]$ (as the other inclusion is obvious). We will need the following analog of Lemma 31. The difference is that the ``blow-up'' of the time bound in $\mathrm{pK}^t$ here is linear in $t$, while in Lemma 31, this blow-up is polynomial in $t$.
\paragraph{Lemma 46 (Fine-Grained Computational Depth).} Let $x\in\{0,1\}^n$ and $r\in\{0,1\}^*$ be arbitrary strings, $f$ any function, and $\tau_n(t):=t\cdot\operatorname{poly}(n)$. There exists a constant $a>0$ and a time bound $t$ such that $f(n)\leq t\leq f(n)\cdot2^{a\sqrt{n\log n}}$ and
\[\mathrm{pK}^t(x\mid r)-\mathrm{pK}^{\tau_n(t)}(x\mid r)\leq O(\sqrt{n\log n}).\]
\paragraph{Proof.} Fix $x\in\{0,1\}^n$. Let $\tau_n(t):=t\cdot n^c$. Define $t_0:=f(n)$, and for $i\geq1$, $t_i:=\tau_n(t_{i-1})$. Consider the following telescoping sum:
\begin{align*}
\mathrm{pK}^{t_0}(x\mid r)-\mathrm{pK}^{t_I}(x\mid r)
&=(\mathrm{pK}^{t_0}(x\mid r)-\mathrm{pK}^{t_1}(x\mid r))\\
&\quad+(\mathrm{pK}^{t_1}(x\mid r)-\mathrm{pK}^{t_2}(x\mid r))\\
&\quad+\cdots+(\mathrm{pK}^{t_{I-1}}(x\mid r)-\mathrm{pK}^{t_I}(x\mid r)).
\end{align*}
Note that $\mathrm{pK}^{f(n)}(x\mid r)\leq n+d$, for some universal constant $d\geq0$; hence, the above sum is at most $n+d$. By averaging, there is some index $1\leq i^*\leq I$ such that
\[\mathrm{pK}^{t_{i^*-1}}(x\mid r)-\mathrm{pK}^{t_{i^*}(n)}(x\mid r)\leq(n+d)/I,\]
which is at most $O(\sqrt{n\log n})$ for $I:=\sqrt{n/\log n}$. Finally, by induction, for every $i\geq0$, $t_i\leq f(n)\cdot2^{c\cdot i\cdot\log n}$. The bound $t_{i^*-1}\leq f(n)\cdot2^{O(\sqrt{n\cdot\log n})}$ follows.\hfill$\triangleleft$
We will also need the following fine-grained version of Lemma 26.
\paragraph{Lemma 47 (Fine-Grained Symmetry of Information for $\mathrm{pK}^t$).} If $\mathrm{NQL}\times\mathrm{QLSamp}\subseteq\mathrm{AvgBPQL}$, then there exists a polynomial $p$ such that for all sufficiently large $n,m\in\mathbb N$, all $t\in\mathbb N$ such that $t\geq2^{(n\cdot m)^\varepsilon}$ for some $\varepsilon>0$, and all time-constructible $0\leq\tau(n,m)\leq t$, the following holds: for every $x\in\{0,1\}^n$ and every family $\{y_r\in\{0,1\}^m\mid r\in\{0,1\}^{\tau(n,m)}\}$, with probability at least $9/10$ over $r\in\{0,1\}^{\tau(n,m)}$,
\[\mathrm{pK}^t(x,y_r\mid r)>\mathrm{pK}^{\widetilde O(t)}(x\mid r)+\mathrm{pK}^{\widetilde O(t)}(y_r\mid x,r)-\log p(t).\]
\paragraph{Proof Sketch.} The proof is adapted from that of Lemma 26. Define
\begin{multline*}
L:=\{(u,v,w,w',1^t,1^s,1^n,1^m)\mid\exists M\in\{0,1\}^s,M(w,w')\text{ prints }uv\text{ within }2t\text{ steps, where}\\
|w'|=\tau(n,m)\leq t,|w|=2t,\text{ and }s=|u|+|v|-10\}.
\end{multline*}
A distribution $D$ will be given parameters $\langle nk+k,mk'+k',4t,s,n,m\rangle=:N$. It will be defined to randomly sample $u\sim U_{nk+k}$, $v\sim U_{mk'+k'}$, $w\sim U_{2t}$, and $w'\sim U_{\tau(n,m)}$, and then output $(u,v,w,w',1^t,1^s,1^n,1^m)$.
The key observation is that $L$ can be solved in NQL, so using our fine-grained average-case easiness assumption, we obtain a randomized heuristic algorithm $B$ that solves the distributional problem $(L,D)$ in quasilinear time. In particular, since $t\geq2^{(n\cdot m)^\varepsilon}$ for some $\varepsilon>0$, $N\leq\widetilde O(t)$. This ensures that $B$ runs in time $\widetilde O(t)$ (see Proposition 16 and Definition 5).
The rest of the proof follows that of Lemma 26, except when $B$ is used later on as a distinguisher, where we invoke the reconstruction lemma (Lemma 22), we get time bounds of $\widetilde O(t)$ instead of $\operatorname{poly}(t)$.\hfill$\triangleleft$
Using similar ideas, we can get an analog of Lemma 33 for a special case of $\ell=1$. We will use the following definitions, similar to Definition 32, which respectively allow for larger time-bounds and impose a separate length restriction on witnesses.
\paragraph{Definition 48 ($\mathrm{BP}_\delta\circ\mathrm{NTIME}$ and $\mathrm{BP}_\delta\circ\mathrm{NTIMEGUESS}$).} A language $L$ is in the class $\mathrm{BP}_\delta\circ\mathrm{NTIME}[\tau]$ for a time bound $\tau$ if there is a formula
\[\varphi(x)=\mathrm{BP}_r\ \exists y\ R(x,y,r)\]
for a predicate $R$ running in time $\tau$ and variables $y,r$ of length at most $\tau$, such that, for all $x\in\{0,1\}^n$,
\begin{align*}
x\in L&\Longrightarrow\Pr_r[\exists y\ R(x,y,r)]\geq\delta,\text{ and}\\
x\notin L&\Longrightarrow\Pr_r[\forall y\ \neg R(x,y,r)]\geq\delta.
\end{align*}
The class $\mathrm{BP}_\delta\circ\mathrm{NTIMEGUESS}[\tau,\ell]$ is defined the same way as above, with an extra condition that $|y|\leq\ell$ (but the randomness $r$ may still be of length at most $\tau$).
\paragraph{Lemma 49 (Fine-Grained Oracle Elimination).} Suppose $\mathrm{NQL}\times\mathrm{QLSamp}\subseteq\mathrm{AvgBPQL}$. Let $\tau(n)=2^{O(\sqrt{n\log n})}$ be some time-constructible function and $c>0$ a constant. For an arbitrary $L\in\mathrm{BP}_\delta\circ\mathrm{NTIMEGUESS}[\tau(n),n^c]$, let $x\in L_n$ be sufficiently large, let $r\in\{0,1\}^{\tau(n)}$ be any random string under the BP quantifier such that there exists a witness for $x\in L$, and let $y_{x,r}\in\{0,1\}^{n^c}$ be the lexicographically first such $L$-witness for $x$ and randomness $r$.
There exist polynomials $q$ and $q_0$ such that, for all sufficiently large $n\in\mathbb N$, all $x\in L_n$, and all $t\geq q_0(\tau(n))$, we have with probability at least $\delta-1/10$ over $r\in\{0,1\}^{\tau(n)}$ that an $L$-witness $y_{x,r}$ exists, and
\[\mathrm{pK}^{\widetilde O(t)}(x,y_{x,r}\mid r)\leq\mathrm{pK}^t(x\mid r)+\log q(t).\]
\paragraph{Proof Sketch.} The proof is adapted from that of Lemma 33. Define
\begin{multline*}
L':=\{(\mathrm{DP}_k(x,y;z),w,w',1^t,1^k,1^s,1^n)\mid\exists M\in\{0,1\}^s,\text{ such that }M^{\mathrm{SAT}}(w,w')\text{ prints}\\
(x,y)\in\{0,1\}^{n+n^c}\text{ within }2t\text{ steps, where }|w'|=\tau(n)\leq t,|w|=2t,\text{ and }s=k-10\}.
\end{multline*}
Since SAT is NQL-complete under many-one QL reduction, we have, for $x\in\{0,1\}^n$, $y\in\{0,1\}^{n^c}$, $w\in\{0,1\}^{2t}$, and $w'\in\{0,1\}^{\tau(n)}$ that $L'\in\Sigma_2^{\mathrm{QL}}$.
Again, using our fine-grained average-case easiness assumption, we get a randomized heuristic algorithm $B$ that solves $L'$ in time $\widetilde O(t)$. Also note that given $x$ and good randomness $r\in\{0,1\}^{\tau(n)}$, we can compute $y_{x,r}$ via search-to-decision reduction using a SAT oracle, which takes time $\operatorname{poly}(\tau)=:q_0(\tau)$. Therefore, similar to Equation (19) in the original proof, we get
\[\mathrm{pK}^{2t,\mathrm{SAT}}(x,y_{x,r}\mid r)\leq\mathrm{pK}^t(x\mid r)+d\log t\tag{30}\]
for some constant $d>0$, provided $t\geq q_0(\tau)$. The rest of the proof is the same as that of Lemma 33, except when $B$ is used later in the proof as a distinguisher, where we invoke the reconstruction lemma (Lemma 22), we get a time bound of $\widetilde O(t)$ instead of $\operatorname{poly}(t)$.\hfill$\triangleleft$
We will also use the following simple generalization of (non-fine-grained) oracle elimination, Lemma 33, to higher time bounds.
\paragraph{Lemma 50.} Suppose $\mathrm{Dist}\Sigma_2^{\mathrm P}\subseteq\mathrm{AvgBPP}$. For an arbitrary $L\in\mathrm{NTIME}[\tau(n)]$, for some $\tau(n)=2^{O(\sqrt{n\log n})}$, let $x\in L_n$ be sufficiently large, and let $y_x\in\{0,1\}^{\tau(n)}$ be the lexicographically first $L$-witness for $x$. There exist polynomials $q$ and $q_0$ such that, for all $t\geq q_0(\tau(n))$, we have
\[\mathrm{pK}^{q(t)}(x,y_x)\leq\mathrm{pK}^t(x)+\log q(t).\]
\paragraph{Proof Sketch.} The proof follows that of the ``Moreover'' statement of Lemma 33, except the search-to-decision computation of $y_x$, where $|y_x|=\tau(n)$, now takes time at least $q_0(\tau(n))$ for some polynomial $q_0$, yielding
\[\mathrm{pK}^{2t,\mathrm{NP}}(x,y_{x,r}\mid r)\leq\mathrm{pK}^t(x\mid r)+d\log n\]
for $t\geq q_0(\tau(n))$ in Equation (19).\hfill$\triangleleft$
We are now ready to show Theorem 45.
\paragraph{Proof of Theorem 45.} Assume $\Sigma_2^{\mathrm{QL}}\times\mathrm{QLSamp}\subseteq\mathrm{AvgBPQL}$. Let $\tau=2^{a\sqrt{n\log n}}$ for some constant $a>0$, and let $L\in\mathrm{AMTIME}[\tau]$ be arbitrary. That is, for every $x\in\{0,1\}^n$,
\begin{align*}
x\in L&\Longrightarrow\mathrm{BP}_{(1-2^{-n})}r\ \exists y\ V(x,y,r),\\
x\notin L&\Longrightarrow\mathrm{BP}_{(1-2^{-n})}r\ \forall y\ \neg V(x,y,r),
\end{align*}
for a polytime verifier $V(x,y,r)$, where $|r|,|y|=\tau$. Let $x\in L_n$ and good randomness $r\in\{0,1\}^\tau$ be given, with lexicographically first witness $y_{x,r}\in\{0,1\}^\tau$.
Note that $\Sigma_2^{\mathrm{QL}}\times\mathrm{QLSamp}\subseteq\mathrm{AvgBPQL}$ implies $\mathrm{Dist}\Sigma_2^{\mathrm P}\subseteq\mathrm{AvgBPP}$ by a simple padding argument. Combining Lemma 50 and Lemma 26 as in the proof of Lemma 34, for some constant $b>0$, with probability at least $1-2^{-n}-1/5$ over $r\in\{0,1\}^\tau$,
\[\mathrm{pK}^{\tau^b}(y_x\mid x,r)\leq(\mathrm{pK}^\tau(x\mid r)-\mathrm{pK}^{\tau^b}(x\mid r))+O(\log\tau),\tag{31}\]
from which it follows that $\mathrm{pK}^{\tau^b}(y_x\mid x,r)\leq2n$ by the trivial upper bound on $\mathrm{pK}^\tau(x\mid r)$. In other words, by definition of pK,
\[\Pr_{r'\sim\{0,1\}^{\tau^b}}[\exists M\in\{0,1\}^{2n},M(x,r,r')\text{ outputs }y_x\text{ within time }\tau^b]\geq2/3.\tag{32}\]
Define a verifier $V'(x,z,r,r')$, which does the following:
\begin{itemize}
\item Simulate the machine described by $z$ on input $(x,r,r')$ for $\tau^b$ steps to produce some output $y_z$, and then return the output of $V(x,y_z,r)$.
\end{itemize}
Note that $V'$ runs in time $\tau^c$ for some constant $c>0$. For $x\in L$, there exists $z\in\{0,1\}^{2n}$ such that $V'(x,z,r,r')$ so long as Equation (31) and Equation (32) both hold, which occurs with probability at least $(2/3)(1-2^{-n}-1/5)=:\delta>1/2$ over $r,r'$. For $x\notin L$, $V'(x,z,r,r')$ rejects unless there is an $L$-witness $y_{x,r}$ under the original verifier $V$, which occurs with probability at most $2^{-n}$. This leads to the following characterization of $L$:
\begin{align*}
x\in L&\Longrightarrow\mathrm{BP}_{(\delta)}r,r'\ \exists z\ V'(x,z,r.r'),\\
x\notin L&\Longrightarrow\mathrm{BP}_{(1-2^{-n})}r,r'\ \forall z\ \neg V'(x,z,r,r'),
\end{align*}
where $|r|=\tau$, $|r'|=\tau^b$, and $|z|=2n$, which implies that $L\in\mathrm{BP}_\delta\circ\mathrm{NTIMEGUESS}[\tau^c,2n]$.
Combining Lemma 47 and Lemma 49 as in the proof of Lemma 34, we have that with probability at least $\delta-1/5>1/4$ over $(r,r')\in\{0,1\}^{\tau+\tau^b}$ that for all $t\geq p'_0(\tau^c)$, for some polynomials $p'$ and $p'_0$,
\[\mathrm{pK}^{t\cdot p'(n)}(z_{x,r,r'}\mid x,r,r')\leq(\mathrm{pK}^t(x\mid r,r')-\mathrm{pK}^{t\cdot p'(n)}(x\mid r,r'))+O(\log t)\tag{33}\]
where $z_{x,r,r'}$ is such that $V'(x,z_{x,r,r'},r,r')$ accepts. Lemma 46 then implies the existence of some $t\in\mathbb N$ such that $p'_0(\tau^c)\leq t\leq2^{O(\sqrt{n\log n})}$ and $\mathrm{pK}^t(x\mid r,r')-\mathrm{pK}^{t\cdot p'(n)}(x\mid r,r')\leq O(\sqrt{n\log n})$. So by Equation (33), for some constant $d>0$,
\[\mathrm{pK}^{\tau^d}(z_{x,r}\mid x,r,r')\leq d\sqrt{n\log n}.\]
This suggests the following algorithm to decide $L$:
\begin{itemize}
\item Given $x\in\{0,1\}^n$, sample uniformly random strings $r\sim\{0,1\}^\tau$, $r'\in\{0,1\}^{\tau^b}$, and $w\in\{0,1\}^{\tau^d}$. Then exhaustively search over all $d\sqrt{n\log n}$-bit machines $M$, running each $M(x,r,r',w)$ for $\tau^d$ steps to produce a candidate witness $z$, accepting if $V'(x,z,r,r')$ accepts.
\end{itemize}
The runtime $\operatorname{poly}(\tau)=2^{O(\sqrt{n\log n})}$ is clear. For correctness, for every $x\in L_n$, there is at least a $1/4$ fraction of good random strings $r,r'$ such that Equation (33) holds, in which case a witness $z_{x,r,r'}$ will be found by our algorithm with probability at least $2/3$ over $w$ (by definition of pK). So, our algorithm accepts $x\in L_n$ with probability at least $(1/4)(2/3)=1/6$. On the other hand, for $x\notin L$, our algorithm will accept with probability at most $2^{-n}$, as this upper-bounds the probability of $r$ being good for the original verifier $V$. Thus, we have $L\in\mathrm{BPTIME}[2^{O(\sqrt{n\log n})}]$ via standard success amplification.\hfill$\triangleleft$
A proof of the analogous statement for the class $\mathrm{UTIME}[2^{O(\sqrt{n\log n})}]$, Item 1 of Theorem 2,
\[\mathrm{NQL}\times\mathrm{QLSamp}\subseteq\mathrm{AvgBPQL}\Longrightarrow\mathrm{UTIME}[2^{O(\sqrt{n\log n})}]\subseteq\mathrm{RTIME}[2^{O(\sqrt{n\log n})}],\]
is presented in Appendix B. Though its proof requires some additional ideas, it does not need to use an extremely efficient PRG as the one developed in [6].

% Faithful official CCC2022 article16 full AppendixB.
\section*{B Probabilistic Fine-Grained Reduction for $\mathrm{UTIME}[2^{O(\sqrt{n\log n})}]$}

\paragraph{Theorem 67.} $\mathrm{NQL}\times\mathrm{QLSamp}\subseteq\mathrm{AvgBPQL}\Longrightarrow\mathrm{UTIME}[2^{O(\sqrt{n\log n})}]\subseteq\mathrm{RTIME}[2^{O(\sqrt{n\log n})}]$.

We will need the following fine-grained version of Lemma 27.
\paragraph{Lemma 68.} If $\mathrm{NQL}\times\mathrm{QLSamp}\subseteq\mathrm{AvgBPQL}$, then there exists a polynomial $\tau$ such that the following promise problem is in promiseBPP:
\begin{align*}
\Pi_{\mathrm{YES}}&:=\{(x,1^s,1^t)\mid\mathrm{pK}^t(x)\leq s\},\\
\Pi_{\mathrm{NO}}&:=\{(x,1^s,1^t)\mid\mathrm{pK}^{\widetilde O(t)\cdot\tau(|x|)}(x)>s+\log\tau(t)\}.
\end{align*}
\paragraph{Proof.} Define
\[L:=\{(\mathrm{DP}_k(x;z),w,1^s,1^n)\mid\exists M\in\{0,1\}^s,M(w)\text{ outputs }x\in\{0,1\}^n\text{ within }|w|\text{ steps}\},\]
where $c>0$ is some large constant so that $L\in\mathrm{NQL}$. Define a distribution family $\mathcal D:=\{D_{\langle nk+k,t,s,n\rangle}\}$, each member of which does the following: sample $u\sim U_{nk+k}$ and $w\sim U_t$ and then output $(u,w,1^s,1^n)$. (Again, for the simplicity of presentation, we omit the padding, which is of length at most $\widetilde O(t)\cdot\operatorname{poly}(nk)$.) By assumption, $(L,\mathcal D)\in\mathrm{AvgBPQL}$. Let $B$ be a randomized heuristic algorithm for $(L,\mathcal D)$ as described in Lemma 6. Now, define an algorithm $B'$:
\begin{itemize}
\item On input $(x,1^s,1^t)$ with $x\in\{0,1\}^n$, set $k=s+10$, sample $z\sim U_{nk}$ and $w\sim U_t$, and then output $B(\mathrm{DP}_k(x;z),w,1^s,1^n)$.
\end{itemize}
Note that $B'$ runs in time $t_D:=\widetilde O(t)\cdot p(n)$ for some polynomial $p$. Below, we show that $B'$ solves the mentioned promise problem correctly with high probability in the worst case.
First, consider the case that $(x,1^s,1^t)\in\Pi_{\mathrm{YES}}$. By the definitions of $L$ and pK, for any choice of $z\in\{0,1\}^{nk}$,
\[\Pr_w[(\mathrm{DP}_k(x;z),w,1^s,1^n)\in L]\geq2/3.\]
The definition of $B$ then implies that
\[\Pr_{w,z,r_B}[B(\mathrm{DP}_k(x;z),w,1^s,1^n)=1]>1/2,\]
and so
\[\Pr_{r_{B'}}(B'(x,1^s,1^t)=1)>1/2,\tag{42}\]
where $r_B$ denotes the internal randomness of $B$, and $r_{B'}=(w,z,r_B)$ that of $B'$.
Now consider the case that $(x,1^s,1^t)\in\Pi_{\mathrm{NO}}$. For a contradiction, suppose that
\[\Pr_{w,z,r_B}[B(\mathrm{DP}_k(x;z),w,1^s,1^n)=1]>1/3.\tag{43}\]
Recall $k=s+10$. By a counting argument, for randomly selected $u$ and $w$,
\[\Pr_{u,w}[(u,w,1^s,1^n)\in L]\leq\frac{2^s\cdot2^{nk}\cdot2^t}{2^{nk+k+t}}=1/10.\]
Then by definition of $B$,
\[\Pr_{u,w,r_B}[B(u,w,1^s,1^n)=1]\leq1/9.\tag{44}\]
Comparing Equations (43) and (44), we see that $B(-,U_t,1^s,1^n)$ $(2/9)$-distinguishes the distribution $\mathrm{DP}_k(x;U_{nk})$ from uniform. Lemma 22 implies that
\[\mathrm{pK}^{\widetilde O(t)\cdot\tau_0(n)}(x)\leq k+\log\tau_0(t)=s+10+\log\tau_0(t),\]
for some polynomial $\tau_0$ that depends on the polynomials $p$ and $p_{\mathrm{DP}}$ from Lemma 22. This means that $(x,1^s,1^t)$ is not in $\Pi_{\mathrm{NO}}$, by setting $\tau$ properly. This gives the desired contradiction. By definition of $B'$, we have that
\[\Pr_{r_{B'}}(B'(x,1^s,1^t)=1)\leq1/3.\tag{45}\]
By Equations (13) and (16), $B'$ yields a promiseBPP algorithm for $(\Pi_{\mathrm{YES}},\Pi_{\mathrm{NO}})$ via standard error reduction techniques.\hfill$\triangleleft$
The following is a simple corollary of Lemma 47 and Lemma 20.
\paragraph{Lemma 69 (Fine-Grained Weak Symmetry of Information for $\mathrm{pK}^t$).} If $\mathrm{NQL}\times\mathrm{QLSamp}\subseteq\mathrm{AvgBPQL}$, then there exist polynomials $p_0$ and $p$ such that for every sufficiently large $x\in\{0,1\}^n$, $m\in\mathbb N$ and $p_0(n,m)\leq t\leq2^{n+m}$,
\[\Pr_{u\sim\{0,1\}^m}[\mathrm{pK}^t(x\circ u)>\mathrm{pK}^{t\cdot p(nm)}(x)+m-\log p(t)]\geq0.99.\]
We will use the following black-box hitting set generator whose outputs can be small relative to the ``hard string'' used in its construction.
\paragraph{Lemma 70 ([10, Theorem 4.3]).} For any $T,m\in\mathbb N$ with $m\leq2T$, there exists a function $H_{T,m}:\{0,1\}^T\times\{0,1\}^d\to\{0,1\}^m$ and a deterministic procedure $\mathrm{Recon}^{(-)}:\{0,1\}^a\to\{0,1\}^T$, where $d=O(\log T+\log^3m)$ and $a=2m+O(\log T+\log^3m)$ such that, for any $x\in\{0,1\}^T$ and any function $D:\{0,1\}^m\to\{0,1\}$ that $0.1$-avoids $H_{T,m}(x,-)$, there exists advice $\alpha\in\{0,1\}^a$ such that $\mathrm{Recon}^D(\alpha)=x$. Moreover, $H_{T,m}$ can be computed in time $\operatorname{poly}(T)$ and $\mathrm{Recon}^D$ can be computed in time $\operatorname{poly}(T)$ with oracle access to $D$.
We are now ready to show Theorem 67.
\paragraph{Proof of Theorem 67.} Let $L\in\mathrm{UTIME}[T(n)]$ with verifier $V$, where $T(n)=2^{O(\sqrt{n\log n})}$. Fix an input $x\in\{0,1\}^n$. Let $y_x\in\{0,1\}^{T(n)}$ be the unique $L$-witness for $x$. Let $B$ be a promiseBPP algorithm that solves the promise problem $(\Pi_{\mathrm{YES}},\Pi_{\mathrm{NO}})$ from Lemma 68.
Let $t$ be any number such that $q_0(T(n))\leq t\leq2^{n/2}$, where $q_0$ is some polynomial specified later. Let $H:=H_{T(n),m}$ be the hitting set generator from Lemma 70, where $m:=m_x(t)\leq O(n)$ is defined later, and let $d:=O(\log T(n)+\log^3m)$ be the seed length of $H$.
Using $H$ and $B$, we will argue that the time-bounded pK complexity of $y_x$ is small given $x$. First note the following.
\paragraph{Claim 71.} There exists a polynomial $q$ such that for every $q_0(T(n))\leq t\leq2^{n/2}$ and every $z_H\in\{0,1\}^d$,
\[\mathrm{pK}^{t\cdot q(nm)}(x\circ H(y_x;z_H))\leq\mathrm{pK}^t(x)+d+\log q(t).\]
\paragraph{Proof of Claim 71.} Let $L$ be the following language
\begin{multline*}
L':=\{(\mathrm{DP}_k(x\circ H(y;z');z),w01^{T(n)^c},1^s,1^n,1^m)\mid\exists M\in\{0,1\}^s,\\
M(w)\text{ outputs }x\in\{0,1\}^n\text{ within }|w|\text{ steps, and }V(x,y)=1\},
\end{multline*}
where $c>0$ is some constant so that $L'\in\mathrm{NQL}$. Define a distribution family $\mathcal D:=\{D_{\langle(n+m)k+k,t+T(n)^c+1,s,n,m\rangle}\}$, each member of which samples $u\sim U_{(n+m)k+k}$, $w\sim U_t$ and outputs $(u,w01^{T(n)^c},1^s,1^n,1^m)$.
By assumption, $(L',\mathcal D)\in\mathrm{AvgBPQL}$. Let $B'$ be a randomized heuristic algorithm for $(L',\mathcal D)$ as described in Lemma 6.
Define $s:=\mathrm{pK}^t(x)$. By definition of $\mathrm{pK}^t$ and $L'$, for any fixed $z\in\{0,1\}^{(n+m)k}$,
\[\Pr_w[(\mathrm{DP}_k(x\circ H(y_x;z_H);z),w01^{T(n)^c},1^s,1^n,1^m)\in L']\geq2/3.\]
Then by definition of $B'$, it is easy to see that
\[\Pr_{z,w,r_{B'}}[B'(\mathrm{DP}_k(x\circ H(y_x;z_H);z),w01^{T(n)^c},1^s,1^n,1^m)=1]\geq1/2.\tag{46}\]
On the other hand, by a counting argument it is easy to see that for $u$ and $w$ selected uniformly at random, we have
\[\Pr_{u,w}[(u,w01^{T(n)^c},1^s,1^n,1^m)\in L']\leq\frac{2^s\cdot2^t\cdot2^d\cdot2^{(n+m)k}}{2^{|u|}\cdot2^{|w|}}\leq1/n,\]
where the last inequality holds by setting $k:=s+d+\log n$. By the definition of $B'$, we have
\[\Pr_{u,w,r_{B'}}[B'(u,w01^{T(n)^c},1^s,1^n,1^m)=1]\leq o(1).\tag{47}\]
Given Equations (46) and (47), it is clear that $D(-):=B'(-,U_t01^{T(n)^c},1^s,1^n,1^m)$ is a randomized distinguisher for $\mathrm{DP}_k(x\circ H(y_x;z_H);U_{(n+m)k})$. Note that provided that $q_0$ is a large enough polynomial, $D$ needs only $O(\log t)$ bits as advice, and then it runs in time
\[t_D:=\widetilde O(t+T(n)^c)\cdot\operatorname{poly}(nm)\leq\widetilde O(t)\cdot\operatorname{poly}(nm).\]
Lemma 22 implies that
\[\mathrm{pK}^{\widetilde O(t)\cdot q(nm)}(x\circ H(y_x;z_H))\leq k+O(\log t)\leq\mathrm{pK}^t(x)+d+\log q(t),\]
for some large polynomial $q$. This completes the proof of Claim 71.\hfill$\triangleleft$
By Lemma 69, there exist polynomial $p_0$ and $p$ such that for every $p_0(n,m)\leq t_1\leq2^{n+m}$
\[\Pr_{u\sim\{0,1\}^m}[\mathrm{pK}^{t_1}(x\circ u)>\mathrm{pK}^{t_1\cdot p(nm)}(x)+m-\log p(t_1)]\geq0.99.\tag{48}\]
We set $t_1:=\widetilde O(t\cdot q(nm))\cdot\tau(n+m)$, where $\tau$ is the function from Lemma 68. Note that we can make the polynomial $q_0$ to be large enough so that $t_1$ satisfies the condition for which Equation (48) holds. Now we choose $m$ so that
\[\mathrm{pK}^{t_1\cdot p(nm)}(x)+m-\log p(t_1)\geq s+\log\tau(t\cdot q(nm)).\tag{49}\]
That is, we set
\[m:=(\mathrm{pK}^t(x)-\mathrm{pK}^{t\cdot n^c}(x))+c\cdot\log t,\]
where $c>0$ is some large constant, and we use the fact that $t\leq2^{n/2}$ to simplify $t_1$. Define the following (probabilistic) function:
\[D(-):=\neg B(x\circ-,1^s,1^{t\cdot q(nm)}).\]
Observe the following properties of $D$.
\begin{itemize}
\item For every $z\in\{0,1\}^d$, $\Pr_D[D(H(y_x;z))=0]\geq2/3$. This follows from Claim 71 and the correctness of $B$ on solving the promise problem $(\Pi_{\mathrm{YES}},\Pi_{\mathrm{NO}})$ from Lemma 68.
\item For at least $0.99$ fraction of $\beta\in\{0,1\}^m$, $\Pr_D[D(\beta)=1]\geq2/3$. This follows from Equations (48) and (49) (and the correctness of $B$).
\item $D$ runs in time $\operatorname{poly}(t)$.
\end{itemize}
At this point, we can see that $D$ is a randomized function that distinguishes $H(y_x;U_d)$ from $U_m$, and by randomly fixing its internal randomness we get a deterministic distinguisher with good probability. However, in order to utilize the reconstruction procedure in Lemma 70 to recover $y_x$, we need a deterministic function that avoids $H$. By amplifying the success probability of $D$ on the ``promised'' inputs, we can show that randomly fixing the internal randomness of $D$ gives a function that avoids $H$ with high probability. More specifically, using standard error reduction techniques, we can make both the probability in the first two items above at least $1-1/(4\cdot2^m)$, which only causes a multiplicative overhead of $\operatorname{poly}(m)$ in the running time. By abusing notation, we use the same $D$ to denote this new function with small error. Now consider the set $S$ of inputs for which $D$ has this small error. That is,
\[S:=\{H(y_x;z)\mid z\in\{0,1\}^d\}\cup\{\beta\in\{0,1\}^m\mid\Pr_D[D(\beta)=1]\geq1-1/(4\cdot2^m)\}.\]
Note that $|S|\leq2^m$. For $\beta\in S$, let us say that the correct answer of $\beta$ is $0$ (resp. $1$) if $\beta$ is in the first (resp. second) subset in the definition of $S$. By a union bound, we have
\[\Pr_{r_D\sim\{0,1\}^{\operatorname{poly}(t)}}[\exists\beta\in S\text{ such that }D(\beta;r_D)\text{ is not correct}]\leq|S|\cdot\frac1{4\cdot2^m}\leq1/4,\tag{50}\]
where $r_D$ above denotes the internal randomness of $D$. This means with probability at least $2/3$ over a uniformly random $r_D$, $D(-;r_D)$ is a deterministic function that is correct on all the inputs in $S$. In particular, such a function $0.1$-avoids $H$. Let us say that $r_D$ is good if this is true. Let $\mathrm{Recon}^{(-)}:\{0,1\}^a\to\{0,1\}^{T(n)}$ be the reconstruction procedure in Lemma 70, where $a:=2m+O(\log T+\log^3m)$. We have
\begin{align*}
&\Pr_{r_D\sim\{0,1\}^{\operatorname{poly}(t)}}[\exists\alpha\in\{0,1\}^a\text{ such that }\mathrm{Recon}^{D(-;r_D)}(\alpha)=x]\\
&\geq\Pr_{r_D}[\exists\alpha\in\{0,1\}^a\text{ such that }\mathrm{Recon}^{D(-,r_D)}(\alpha)=x\mid r_D\text{ is good}]\cdot\Pr_{r_D}[r_D\text{ is good}]\\
&\geq2/3.\qquad\text{(Lemma70 and Equation50)}
\end{align*}
Note that given $x$, $D$ can be constructed using the integers $s,t,n,m$ and the code for $B$, which can be encoded using $O(\log t)$ bits. Then the above implies
\[\mathrm{pK}^{\operatorname{poly}(t\cdot T)}(y_x\mid x)\leq a+O(\log t)\leq2m+O(\log t+\log^3m).\tag{51}\]
Note that the above holds for every $q_0(T(n))\leq t\leq2^{n/2}$. Also, recall that $m=\mathrm{pK}^t(x)-\mathrm{pK}^{t\cdot n^c}(x)+c\cdot\log t$ for some constant $c>0$. By Lemma 46, there exists some $t^*$ such that
\[q_0(T(n))\leq t^*\leq q_0(T(n))\cdot2^{c\cdot\sqrt{n\log n}}\leq2^{O(\sqrt{n\log n})},\]
and
\[\mathrm{pK}^{t^*}(x)-\mathrm{pK}^{t^*\cdot n^c}(x)\leq O(\sqrt{n\log n}).\]
Then for such $t^*$, we have $m\leq O(\sqrt{n\log n})$. Therefore, plugging this $t^*$ into Equation (51) for $t$, we have that there exists some constant $d>0$ such that
\[\mathrm{pK}^{2^{d\cdot\sqrt{n\log n}}}(y_x\mid x)\leq d\cdot\sqrt{n\log n}.\]
This suggests the following probabilistic algorithm $A$ for solving $L$:
\begin{itemize}
\item On input $x\in\{0,1\}^n$, $A$ samples $w\sim\{0,1\}^{d\cdot\sqrt{n\log n}}$. It then exhaustively searches over all $M\in\{0,1\}^{d\cdot\sqrt{n\log n}}$, running $M(w,x)$ for $2^{d\cdot\sqrt{n\log n}}$ steps to produce some output $y$. $A$ accepts iff a $y$ is obtained such that $V(x,y)=1$.
\end{itemize}
This completes the proof.\hfill$\triangleleft$

\begin{thebibliography}{45}
\bibitem[1]{ref1} Luis Filipe Coelho Antunes and Lance Fortnow. Worst-case running times for average-case algorithms. In Proceedings of the 24th Annual IEEE Conference on Computational Complexity, CCC 2009, Paris, France, 15-18 July 2009, pages 298–303. IEEE Computer Society, 2009. doi:10.1109/CCC.2009.12.
\bibitem[2]{ref2} Shai Ben-David, Benny Chor, Oded Goldreich, and Michael Luby. On the theory of average case complexity. J. Comput. Syst. Sci., 44(2):193–219, 1992. doi:10.1016/0022-0000(92)90019-F.
\bibitem[3]{ref3} Andrej Bogdanov and Luca Trevisan. Average-case complexity. Found. Trends Theor. Comput. Sci., 2(1), 2006. doi:10.1561/0400000004.
\bibitem[4]{ref4} Andrej Bogdanov and Luca Trevisan. On worst-case to average-case reductions for NP problems. SIAM J. Comput., 36(4):1119–1159, 2006. doi:10.1137/S0097539705446974.
\bibitem[5]{ref5} Harry Buhrman, Lance Fortnow, and Aduri Pavan. Some results on derandomization. Theory Comput. Syst., 38(2):211–227, 2005. doi:10.1007/s00224-004-1194-y.
\bibitem[6]{ref6} Lijie Chen, Shuichi Hirahara, and Neekon Vafa. Average-case hardness of NP and PH from worst-case fine-grained assumptions. In Mark Braverman, editor, 13th Innovations in Theoretical Computer Science Conference, ITCS 2022, January 31 - February 3, 2022, Berkeley, CA, USA, volume 215 of LIPIcs, pages 45:1–45:16. Schloss Dagstuhl - Leibniz-Zentrum für Informatik, 2022. doi:10.4230/LIPIcs.ITCS.2022.45.
\bibitem[7]{ref7} Halley Goldberg and Valentine Kabanets. A simpler proof of the worst-case to average-case reduction for polynomial hierarchy via symmetry of information. Electron. Colloquium Comput. Complex., 7:1–14, 2022. URL: https://eccc.weizmann.ac.il/report/2022/007/.
\bibitem[8]{ref8} Oded Goldreich. The Foundations of Cryptography - Volume 1: Basic Techniques. Cambridge University Press, 2001. doi:10.1017/CBO9780511546891.
\bibitem[9]{ref9} Oded Goldreich and Leonid A. Levin. A hard-core predicate for all one-way functions. In Proceedings of the Twenty-First Annual ACM Symposium on Theory of Computing, STOC ’89, pages 25–32, New York, NY, USA, 1989. Association for Computing Machinery. doi:10.1145/73007.73010.
\bibitem[10]{ref10} Shuichi Hirahara. Characterizing average-case complexity of PH by worst-case meta-complexity. In Sandy Irani, editor, 61st IEEE Annual Symposium on Foundations of Computer Science, FOCS 2020, Durham, NC, USA, November 16-19, 2020, pages 50–60. IEEE, 2020. doi:10.1109/FOCS46700.2020.00014.
\bibitem[11]{ref11} Shuichi Hirahara. Average-case hardness of NP from exponential worst-case hardness assump- tions. In Samir Khuller and Virginia Vassilevska Williams, editors, STOC ’21: 53rd Annual ACM SIGACT Symposium on Theory of Computing, Virtual Event, Italy, June 21-25, 2021, pages 292–302. ACM, 2021. doi:10.1145/3406325.3451065.
\bibitem[12]{ref12} Shuichi Hirahara. Symmetry of information in heuristica. Manuscript, 2021.
\bibitem[13]{ref13} Shuichi Hirahara and Mikito Nanashima. On worst-case learning in relativized heuristica. In Symposium on Foundations of Computer Science (FOCS), 2021.
\bibitem[14]{ref14} Shuichi Hirahara and Rahul Santhanam. Errorless versus error-prone average-case complexity. In Mark Braverman, editor, 13th Innovations in Theoretical Computer Science Conference, ITCS 2022, January 31 - February 3, 2022, Berkeley, CA, USA, volume 215 of LIPIcs, pages 84:1–84:23. Schloss Dagstuhl - Leibniz-Zentrum für Informatik, 2022. doi:10.4230/LIPIcs.ITCS.2022.84.
\bibitem[15]{ref15} Shuichi Hirahara and Rahul Santhanam. Excluding PH pessiland. In Mark Braverman, editor, 13th Innovations in Theoretical Computer Science Conference, ITCS 2022, January 31 - February 3, 2022, Berkeley, CA, USA, volume 215 of LIPIcs, pages 85:1–85:25. Schloss Dagstuhl - Leibniz-Zentrum für Informatik, 2022. doi:10.4230/LIPIcs.ITCS.2022.85.
\bibitem[16]{ref16} Russell Impagliazzo. A personal view of average-case complexity. In Proceedings of the Tenth Annual Structure in Complexity Theory Conference, Minneapolis, Minnesota, USA, June 19-22, 1995, pages 134–147, 1995. doi:10.1109/SCT.1995.514853.
\bibitem[17]{ref17} Russell Impagliazzo. Relativized separations of worst-case and average-case complexities for NP. In Proceedings of the 26th Annual IEEE Conference on Computational Complexity, CCC 2011, San Jose, California, USA, June 8-10, 2011, pages 104–114, 2011. doi:10.1109/CCC.2011.34.
\bibitem[18]{ref18} Russell Impagliazzo and Leonid A. Levin. No better ways to generate hard NP instances than picking uniformly at random. In 31st Annual Symposium on Foundations of Computer Science, St. Louis, Missouri, USA, October 22-24, 1990, Volume II, pages 812–821. IEEE Computer Society, 1990. doi:10.1109/FSCS.1990.89604.
\bibitem[19]{ref19} Russell Impagliazzo and Avi Wigderson. P = BPP if E requires exponential circuits: De- randomizing the XOR lemma. In Frank Thomson Leighton and Peter W. Shor, editors, Proceedings of the Twenty-Ninth Annual ACM Symposium on the Theory of Computing, El Paso, Texas, USA, May 4-6, 1997, pages 220–229. ACM, 1997. doi:10.1145/258533.258590.
\bibitem[20]{ref20} Michael J. Kearns, Robert E. Schapire, and Linda Sellie. Toward efficient agnostic learning. Mach. Learn., 17(2-3):115–141, 1994. doi:10.1007/BF00993468.
\bibitem[21]{ref21} Pravesh K. Kothari and Roi Livni. Improper learning by refuting. In Anna R. Karlin, editor, 9th Innovations in Theoretical Computer Science Conference, ITCS 2018, January 11-14, 2018, Cambridge, MA, USA, volume 94 of LIPIcs, pages 55:1–55:10. Schloss Dagstuhl - Leibniz-Zentrum für Informatik, 2018. doi:10.4230/LIPIcs.ITCS.2018.55.
\bibitem[22]{ref22} Clemens Lautemann. BPP and the polynomial hierarchy. Inf. Process. Lett., 17(4):215–217, 1983. doi:10.1016/0020-0190(83)90044-3.
\bibitem[23]{ref23} Troy Lee. Kolmogorov complexity and formula lower bounds. PhD thesis, University of Amsterdam, 2006.
\bibitem[24]{ref24} Leonid A. Levin. Laws of information conservation (nongrowth) and aspects of the foundation of probability theory. Problemy Peredachi Informatsii, 10(3):30–35, 1974.
\bibitem[25]{ref25} Leonid A. Levin. Average case complete problems. SIAM J. Comput., 15(1):285–286, 1986. doi:10.1137/0215020.
\bibitem[26]{ref26} Ming Li and Paul M. B. Vitányi. An introduction to Kolmogorov complexity and its applications. Springer, 2019.
\bibitem[27]{ref27} Ming Li and Paul M. B. Vitányi. An Introduction to Kolmogorov Complexity and Its Applications, 4th Edition. Texts in Computer Science. Springer, 2019. doi:10.1007/978-3-030-11298-1.
\bibitem[28]{ref28} Yanyi Liu and Rafael Pass. On one-way functions and kolmogorov complexity. In 61st IEEE Annual Symposium on Foundations of Computer Science, FOCS 2020, Durham, NC, USA, November 16-19, 2020, pages 1243–1254, 2020. doi:10.1109/FOCS46700.2020.00118.
\bibitem[29]{ref29} Yanyi Liu and Rafael Pass. On one-way functions from np-complete problems. Electron. Colloquium Comput. Complex., page 59, 2021. URL: https://eccc.weizmann.ac.il/report/2021/059.
\bibitem[30]{ref30} Yanyi Liu and Rafael Pass. On the possibility of basing cryptography on $\mathrm{EXP}\neq\mathrm{BPP}$. In Advances in Cryptology - CRYPTO 2021 - 41st Annual International Cryptology Conference, CRYPTO 2021, Virtual Event, August 16-20, 2021, Proceedings, Part I, pages 11–40, 2021. doi:10.1007/978-3-030-84242-0\_2.
\bibitem[31]{ref31} Luc Longpré and Osamu Watanabe. On symmetry of information and polynomial time invertibility. Inf. Comput., 121(1):14–22, 1995. doi:10.1006/inco.1995.1120.
\bibitem[32]{ref32} Zhenjian Lu and Igor C. Oliveira. An efficient coding theorem via probabilistic representations and its applications. In Nikhil Bansal, Emanuela Merelli, and James Worrell, editors, 48th International Colloquium on Automata, Languages, and Programming, ICALP 2021, July 12-16, 2021, Glasgow, Scotland (Virtual Conference), volume 198 of LIPIcs, pages 94:1–94:20. Schloss Dagstuhl - Leibniz-Zentrum für Informatik, 2021. doi:10.4230/LIPIcs.ICALP.2021.94.
\bibitem[33]{ref33} Zhenjian Lu, Igor C. Oliveira, and Rahul Santhanam. Pseudodeterministic algorithms and the structure of probabilistic time. In Samir Khuller and Virginia Vassilevska Williams, editors, STOC ’21: 53rd Annual ACM SIGACT Symposium on Theory of Computing, Virtual Event, Italy, June 21-25, 2021, pages 303–316. ACM, 2021. doi:10.1145/3406325.3451085.
\bibitem[34]{ref34} Zhenjian Lu, Igor C. Oliveira, and Marius Zimand. Optimal coding theorems in time-bounded kolmogorov complexity. In 49th International Colloquium on Automata, Languages and Programming, 2022.
\bibitem[35]{ref35} Peter Bro Miltersen, N. V. Vinodchandran, and Osamu Watanabe. Super-polynomial versus half-exponential circuit size in the exponential hierarchy. In Takao Asano, Hiroshi Imai, D. T. Lee, Shin-Ichi Nakano, and Takeshi Tokuyama, editors, Computing and Combinatorics, 5th Annual International Conference, COCOON ’99, Tokyo, Japan, July 26-28, 1999, Proceedings, volume 1627 of Lecture Notes in Computer Science, pages 210–220. Springer, 1999. doi:10.1007/3-540-48686-0\_21.
\bibitem[36]{ref36} Ashish V. Naik, Kenneth W. Regan, and D. Sivakumar. On quasilinear-time complexity theory. Theor. Comput. Sci., 148(2):325–349, 1995. doi:10.1016/0304-3975(95)00031-Q.
\bibitem[37]{ref37} Noam Nisan and Avi Wigderson. Hardness vs randomness. J. Comput. Syst. Sci., 49(2):149–167, 1994. doi:10.1016/S0022-0000(05)80043-1.
\bibitem[38]{ref38} Igor C. Oliveira. Randomness and intractability in kolmogorov complexity. In Christel Baier, Ioannis Chatzigiannakis, Paola Flocchini, and Stefano Leonardi, editors, 46th International Colloquium on Automata, Languages, and Programming, ICALP 2019, July 9-12, 2019, Patras, Greece, volume 132 of LIPIcs, pages 32:1–32:14. Schloss Dagstuhl - Leibniz-Zentrum für Informatik, 2019. doi:10.4230/LIPIcs.ICALP.2019.32.
\bibitem[39]{ref39} Igor C. Oliveira and Rahul Santhanam. Pseudodeterministic constructions in subexponential time. In Hamed Hatami, Pierre McKenzie, and Valerie King, editors, Proceedings of the 49th Annual ACM SIGACT Symposium on Theory of Computing, STOC 2017, Montreal, QC, Canada, June 19-23, 2017, pages 665–677. ACM, 2017. doi:10.1145/3055399.3055500.
\bibitem[40]{ref40} Hanlin Ren and Rahul Santhanam. Hardness of KT characterizes parallel cryptography. In 36th Computational Complexity Conference, CCC 2021, July 20-23, 2021, Toronto, Ontario, Canada (Virtual Conference), pages 35:1–35:58, 2021. doi:10.4230/LIPIcs.CCC.2021.35.
\bibitem[41]{ref41} Shai Shalev-Shwartz and Shai Ben-David. Understanding Machine Learning - From Theory to Algorithms. Cambridge University Press, 2014. URL: http://www.cambridge.org/de/academic/subjects/computer-science/pattern-recognition-and-machine-learning/ understanding-machine-learning-theory-algorithms.
\bibitem[42]{ref42} Ronen Shaltiel and Christopher Umans. Simple extractors for all min-entropies and a new pseudorandom generator. J. ACM, 52(2):172–216, 2005. doi:10.1145/1059513.1059516.
\bibitem[43]{ref43} Luca Trevisan and Salil P. Vadhan. Pseudorandomness and average-case complexity via uniform reductions. Computational Complexity, 16(4):331–364, 2007.
\bibitem[44]{ref44} Salil P. Vadhan. On learning vs. refutation. In Satyen Kale and Ohad Shamir, editors, Proceed- ings of the 30th Conference on Learning Theory, COLT 2017, Amsterdam, The Netherlands, 7-10 July 2017, volume 65 of Proceedings of Machine Learning Research, pages 1835–1848. PMLR, 2017. URL: http://proceedings.mlr.press/v65/vadhan17a.html.
\bibitem[45]{ref45} Emanuele Viola. On constructing parallel pseudorandom generators from one-way functions. In 20th Annual IEEE Conference on Computational Complexity (CCC 2005), 11-15 June 2005, San Jose, CA, USA, pages 183–197, 2005. doi:10.1109/CCC.2005.16.
\end{thebibliography}
\end{document}
